\documentclass[msc, a4paper, 12pt]{ppgccufmg} 

\usepackage[english]{babel}
\usepackage[utf8]{inputenc}  
\usepackage[T1]{fontenc}
\usepackage{hyperref} 
\usepackage{rotating}
\usepackage{bm, bbm, dsfont, amsfonts, amsmath}
\usepackage{natbib}
\usepackage{tikz}
\usepackage{multicol,multirow}
\makeatletter
\let\c@lofdepth\relax
\let\c@lotdepth\relax
\makeatother
\usepackage{subfigure}
\usepackage{float}
\usepackage[binary-units]{siunitx}
\usepackage[linesnumbered,lined,boxed]{algorithm2e}
\usepackage{algpseudocode}
\usepackage{enumitem}
\usepackage{undertilde}
\usepackage[labelsep=period]{caption}  

\DeclareSIUnit{\minusalphapercent}{(1-\alpha)\%}
\DeclareSIUnit{\alphapercent}{\alpha\%}

\newcommand{\Mycite}[1]{%
 (\citeauthor{#1},~\citeyear{#1})}

%\selectlanguage{brazil}

\usetikzlibrary{svg.path}

\definecolor{orcidlogocol}{HTML}{A6CE39}
\tikzset{
	orcidlogo/.pic={
		\fill[orcidlogocol] svg{M256,128c0,70.7-57.3,128-128,128C57.3,256,0,198.7,0,128C0,57.3,57.3,0,128,0C198.7,0,256,57.3,256,128z};
		\fill[white] svg{M86.3,186.2H70.9V79.1h15.4v48.4V186.2z}
		svg{M108.9,79.1h41.6c39.6,0,57,28.3,57,53.6c0,27.5-21.5,53.6-56.8,53.6h-41.8V79.1z M124.3,172.4h24.5c34.9,0,42.9-26.5,42.9-39.7c0-21.5-13.7-39.7-43.7-39.7h-23.7V172.4z}
		svg{M88.7,56.8c0,5.5-4.5,10.1-10.1,10.1c-5.6,0-10.1-4.6-10.1-10.1c0-5.6,4.5-10.1,10.1-10.1C84.2,46.7,88.7,51.3,88.7,56.8z};
	}
}

\newcommand\orcidicon[1]{\href{https://orcid.org/#1}{\mbox{\scalerel*{
				\begin{tikzpicture}[yscale=-1,transform shape]
				\pic{orcidlogo};
				\end{tikzpicture}
			}{|}}}}


\renewcommand{\Bbb}{\mathbb}

\newdimen\HilbertLastX
\newdimen\HilbertLastY
\newcounter{HilbertOrder}

\def\DrawToNext#1#2{%
	\advance \HilbertLastX by #1
	\advance \HilbertLastY by #2
	\pgfpathlineto{\pgfqpoint{\HilbertLastX}{\HilbertLastY}}
}

\def\Hilbert[#1,#2,#3,#4,#5,#6,#7,#8] {
	\ifnum\value{HilbertOrder} > 0%
	\addtocounter{HilbertOrder}{-1}
	\Hilbert[#5,#6,#7,#8,#1,#2,#3,#4]
	\DrawToNext {#1} {#2}
	\Hilbert[#1,#2,#3,#4,#5,#6,#7,#8]
	\DrawToNext {#5} {#6}
	\Hilbert[#1,#2,#3,#4,#5,#6,#7,#8]
	\DrawToNext {#3} {#4}
	\Hilbert[#7,#8,#5,#6,#3,#4,#1,#2]
	\addtocounter{HilbertOrder}{1}
	\fi
}

% \hilbert((x,y),order)
\def\hilbert((#1,#2),#3){%
	\advance \HilbertLastX by #1
	\advance \HilbertLastY by #2
	\pgfpathmoveto{\pgfqpoint{\HilbertLastX}{\HilbertLastY}}
	\setcounter{HilbertOrder}{#3}
	\Hilbert[1mm,0mm,-1mm,0mm,0mm,1mm,0mm,-1mm]
	\pgfusepath{stroke}%
}

% correct bad hyphenation here
\hyphenation{op-tical net-works semi-conduc-tor}

\begin{document}

\ppgccufmg{
 title = {Contributions to the Study of Time Series and Images with the Entropy-Complexity Plane},
 authorrev = {Caetano Chagas, Eduarda Tatiane},
 cutter = {D1234p},
 cdu = {519.6*82.10},
 university = {Federal University of Minas Gerais},
 course = {Computer Science},
 address = {Belo Horizonte},
 date = {2021-03},
 advisor = {Heitor Soares Ramos Filho},
 coadvisor = {Alejandro Cesar Frery Orgambide},
 abstract = {Resumo}{./intro_tex/abs_por},
 abstract = [english]{Abstract}{./intro_tex/abs_eng},
 %approval = {img/approvalsheet.eps},
 dedication = {./intro_tex/dedication},
 ack = {./intro_tex/acks},
 keywords={Séries Temporais, Padrões ordinais, Teoria da Informação},
 epigraphtext={Não há alegria mais intensa que a de encontrar um fato que não pode ser entendido nos termos das ideias atualmente aceitas}{Cecilia Payne-Gaposchkin},
  indexkeys={1.~Computação --- Teses. 2.~Redes --- Teses. I.~Orientador.
    II.~Título.},
}


\chapter{Introduction} \label{intro}

\section{Motivation}

%Esta dissertação avança o estado-da-arte\dots
%Histórico: Bandt-Pompe, Aplicações, identificação de lacunas (os tópicos 3.1 a 3.4)

In the last decades we have dealt with a drastic increase in the number of applications coming from data mining, consequently causing an increase in the diversity and volume of information used.
With this, the level of complexity of the investigations, the interdisciplinarity and the number of features necessary to carry out such activities were also increased.
Thus, the study of simple approaches, inexpensive computationally and independent of the type of data for the extraction and characterization of patterns has become fundamental.

One of the fields of study in this area is the application of information theory descriptors.
The information theory developed by Claude Shannon emerged as an interdisciplinary branch, producing countless results both in the theoretical point of view and in the applications in information extraction in signals, covering in its solutions concepts present in Probability, Statistics and Telecommunications.
It currently consists of a powerful tool for the quantification of different levels of order and complexity present in the processes that generate the data.

In the context of non-parametric analysis of time series, a new methodology was proposed by~\cite{PermutationEntropyBandtPompe} for data analysis.
Obtaining a representation of the time series in Bandt-Pompe ordinal patterns, two approaches are currently being applied to acquire a non-parametric probability distribution of the data: the use of frequency histograms or transition graphs.
When opting for this new representation of the data, the resulting distribution becomes less sensitive to outliers and, as it does not depend on any model, it can be applied to a variety of situations.
Despite its simplicity, this method is robust to noise and yields good results in assessing the randomness of a sequence, characterization and classification of signals.



The Bandt-Pompe methodology and its variants have been used successfully in the analysis of many types of dynamics, receiving so far more than \num{1800} citations, according to the Web of Science.
We found works using this approach in several areas of scientific knowledge such as, for example:
distinguishing noise from chaos~\citep{rosso2007distinguishing};
the study of electroencephalography signals using wavelet decomposition~\citep{baravalle2018discriminating,baravalle2018rhythmic};
description of El Niño/Southern Oscillation during the Holocene~\citep{saco2010entropy};
the characterization of household appliances through their energy consumption~\citep{CharacterizationElectricLoadInformationTheoryQuantifiers};
detecting and quantifying stochastic and coherence resonance~\citep{rosso2009detectinga, rosso2009detectingb};
analysis and characterization of economic time series, e.g., stock market, sovereign bonds, credit rating, commodities, and cryptocurrencies~\citep{zunino2010complexity, zunino2012efficiency, bariviera2013efficiency, bariviera2018analysis, Araujo2019permutation};
online signature classification and verification~\citep{ClassificationVerificationOnlineHandwrittenSignatures}.
\cite{InformationTheoryPerspectiveNetworkRobustness} verified the effect of attacks on complex networks by the displacement of points in the $H \times C$ plane.
\citet{CharacterizationVehicleBehaviorInformationTheory} described vehicles' behavior depending on the topology of cities, and
\citet{Chagas2020Characterization} succeeded in expanding the use of such techniques for analyzing textured images corrupted by speckle noise.
\citet{LiborInvisibleHand} identified spurious interventions in the Libor market using the $H\times C$ plane representation.
\citet{echegoyen2020permutation} were able to discriminate between individuals with mild cognitive impairment from those diagnosed with Alzheimer's disease using magnetoencephalography recordings.

Given the probability distribution of the patterns, each time series is then described by a point in the $\mathbbm R^2$ subinterval, the Complexity-Entropy plane.
Two points are well known in this plane:
\begin{enumerate}
    \item \textbf{White noises}, that is, random sequences without any spatial structure, where the entropy presents its maximum value while the complexity is minimal - we can describe statistically the systems present in such locations of the plane by a random variable taken from some distribution; and
    \item \textbf{Deterministic data}, that is, sequences with a periodic structure, where entropy and statistical complexity have their minimum values - we can reconstruct the patterns of such data with only a small portion of patterns.
\end{enumerate}

Through these references, we can characterize time series according to the dynamics of their generation process.
Studies with different applications have managed to obtain relevant results from time series through information on the nature of the data provided by the $H \times C$ plane.
Examples include:
~(i) ~\cite{echegoyen2020permutation} analysis of permutation in magnetoencephalography recordings of individuals suffering from mild cognitive impairment and individuals diagnosed with Alzheimer's disease by trajectories in the $H \times C$ plane,
~(ii) the study of~\cite{InformationTheoryPerspectiveNetworkRobustness} of the effect of attacks on complex networks by shifting their points on the $H \times C$ plane, and
~(iii) the description of vehicle behavior based on the topology of cities carried out by \cite{CharacterizationVehicleBehaviorInformationTheory}.

However, we can see from the examples above that Bandt-Pompe symbolization has a wider range of applications for one-dimensional signals, such as time series.
Thus, it still presents challenges in applications with higher dimensions, such as reducing the loss of spatial correlation between the pixels of an analyzed image.

The aforementioned studies illustrate the importance of the symbolization method in the most diverse areas of knowledge.
This dissertation works in this context and advances in the state of the art with a focus on investigating the main problems still present in the literature: 
~(i) the application of the Bandt-Pompe approach in images, 
%with a specific focus on remote sensing images, 
~(ii) the absence of amplitude weighting methods in transition graphs of ordinal patterns, and 
~(iii) the lack of proposals to build confidence regions.

By proposing an expansion of the applicability of Bandt-Pompe in the image context, this work presents a new perspective of extracting texture characteristics from SAR (Synthetic Aperture Radar) images.
In this way, by linearizing the image samples, we proposed the use of time series analysis tools for a new context: where exists dependence and spatial correlation between the elements.


\section{Goals}

The aim of this work is to develop solutions for the main gaps in the Bandt-Pompe symbolization methodology.
Thus, we propose two solutions: the first proposal to incorporate amplitude information in transition graphs and the first approach to build empirical confidence regions in the Complexity-Entropy plane.
We started the work by presenting a modification of the traditional transition graphs of ordinal patterns, which when applied to the context of analysis of SAR image textures proved to be the first approach to be able to characterize and classify images with results directly comparable to the techniques present in the state of the art.
On the other hand, one of the major problems involving the direct use of causal descriptors, Shannon entropy and statistical complexity, for classification activities is the lack of an appropriate distance metric.
Due to the curvilinear shape of the plane, we verified the need for work focused on building confidence regions and test statistics for such tooling.
Thus, to minimize the impacts caused by the absence of a specific metric, we propose the HC-PCA.
Based on this study, we were able to increase the power of characterization of different classes of time series, where as a use case, we present how the proposed approach manages to determine regions of pure randomness present in the plane.

\section{Contributions}

The main contributions of this dissertation work are:
    
\begin{itemize}
    \item[] \textbf{Weighted Amplitude Transition Graph}: We propose the first approach of transition graphs of weighted ordinal patterns using amplitude information of the analyzed sequences.
    In this way, we were able to reduce the ordinal ambiguity present in transition graphs, thus increasing its characterization power.
    
    \item[] \textbf{Analysis and Classification of SAR Textures using Information Theory}: Through WATG, we propose a new representation of SAR textures, which allows a low-dimension characterization useful for, among other applications, its classification.
    Our approach is robust for rotations and the presence of speckle noise.
    In addition to perfect separation between urban areas, pastures, ocean and forest, the proposed descriptors are interpretable in terms of the degree and structure of the spatial dependence between the observations.
    \item[] \textbf{HC-PCA}: We provide the first contribution in the construction of confidence regions in the Complexity-Entropy Plane, and to measure the similarity of new data sequences with the empirical points we propose the construction of a new test statistic.
    \item[] \textbf{Testing White Noise in the confidence regions}: We present and evaluate a new method of building empirical confidence regions in the Complexity-Entropy plane for analysis of white noise.
    We were able to capture the random behavior of short sequences of PRNGs already analyzed in the literature and we found that in the scenario presented, our technique is robust to correlation structures.
\end{itemize}

\section{Publications}

The first year of the master's degree was dedicated to research an extensive literature review work, where a comprehensive study was review focusing on the classification of SAR textures, alternatives for using Bandt-Pompe symbolization in two-dimensional data, and ordinal pattern weighting approaches by amplitude.
Once have the state-of-the-art, we use the knowledge of several fields of research, such as machine learning, information theory, and graphs to prepare our first proposal, which resulted in some works listed below in chronological order of publication/presentation:
\begin{itemize}
    \item WATG: Incorporating amplitude in ordinal pattern decomposition for time series analysis. Poster presentation at Khipu, Nov 2019, Montevideo, Uruguay. Latin American Meeting In Artificial Intelligence, 2019.
    \item Characterization Of SAR Images With Weighted Amplitude Transition Graphs. 2020 IEEE Latin American GRSS \& ISPRS Remote Sensing Conference (LAGIRS).
    \item Analysis and Classification of SAR Textures using Information Theory. IEEE Journal of Selected Topics in Applied Earth Observations and Remote Sensing (2020).
    \item An exploratory study of the transition graphs application in different resolutions and polarizations of SAR images. 2021 China International SAR Symposium (CISS). (In writing and submission process).
\end{itemize}
During the second year, we plan to focus on developing our proposal for building confidence regions in the Complexity-Entropy plane.
Although they presented unlivable results, detailed studies were carried out involving the use of classical bi-variate analysis, linear regression, and generalized linear models.
After defining the methodology, we focus on writing the article below that is currently in the submission process:
\begin{itemize}
    \item Confidence Regions for Information-Theoretic Descriptors of Time Series. International Statistical Review.
\end{itemize}
A work related to the general area studied was also carried out in collaboration with other researchers:
\begin{itemize}
    \item Detecção de eventos no Twitter através de Grafos de visibilidade natural. III Workshop de Computação Urbana. SBC, 2019.
    \item Supervised Distance Metric learning Encoder with Similarity Space for malware classification through image representation. Computer Networks (Under submission). 
\end{itemize}

\section{Work Organization}

%The rest of this document is organized as follows.
%Chapter~\ref{chapter:BP} introduces concepts and definitions of the Bandt-Pompe symbolization process and the Complexity-Entropy plane.
%In this way, we begin the chapter presenting in the section~\ref{sub:ordinalPatterns} how the ordinal patterns representation is carried out, describing in sections~\ref{sub:rankPermutation} and~\ref{sub:chronological} the permutation modes used and their main differences.
%The concept of transition graphs is described in section~\ref{sub:OPTG}, where we also analyze its main properties.
%Then, in section~\ref{sub:InformationTheory} we discussed about the causal descriptors of Information Theory and how the Complexity-Entropy plane is defined.

This document is organized as follows.
Chapter~\ref{chapter:BP} introduces concepts and definitions of the Bandt-Pompe symbolization process and the Complexity-Entropy plane.
Chapter~\ref{chapter:WATG} provides an overview of ordinal pattern weighting techniques. 
We also propose the first algorithm for incorporating amplitude information into transition graphs, the WATG.
Chapter~\ref{chapter:SARclassification} presents a use case for WATG, providing the experimental results obtained in characterizing and classifying homogeneous textures of SAR images.
Chapter~\ref{chapter:HCPCA} discusses the state of the art analysis of white noise confidence regions in the plane. 
Here we also propose the HC-PCA region of empirical confidence and the construction of a specific statistical test.
Chapter~\ref{chapter:confidenceRegions} presents the main results obtained in our case study with white noise samples.
Finally, in Chapter~\ref{chapter:conclusion}, we provide our final thoughts, a perspective on this work, and directions for future work.
\chapter{Bandt-Pompe Symbolization: Concepts and definitions}\label{chapter:BP}

Dynamic systems describe the relationship in time of a point in a geometric space.
Thus, the study of time series assists in the analysis of the dynamics of its generating processes.
Whether a scalar time series collected through the observation of natural phenomena or obtained through synthetic simulations, its values will be the result of a function formed by its majority of unobserved and measured variables.
Thus, an important research question in data analysis has been:
\begin{quote}
    Given a system and a sample resulting from it whose evolution can be tracked over time, how much information is encoded in this observable about the dynamics of the underlying system and the variables that characterize them?
\end{quote}
Traditionally, the study of time series is carried out under two lines of study, in the domains of time and frequency~\citep{BrockwellDavis91}.
However, both approaches directly use the data resulting from the observational process, which are sensitive to noise and contamination effects.
The approach to the use of non-parametric methods appears in the literature as a way to prevent such effects from compromising the analyzes.

To make good inferences in data analysis studies, it is expected that the proposed approach meets the following requirements:
\begin{itemize}
    \item Be simple, fast and transparent;
    \item Make little or no assumptions about the underlying process; and
    \item Be resilient towards outliers.
\end{itemize}
When analyzing classic probabilistic techniques, we found that they cannot obtain good results without assuming characteristic properties of the data, such as the shape of the probability distribution of the samples.
On the other hand, proposals based on machine learning do not mean to guarantee the required simplicity and transparency, as they do not provide a clear view of the observed characteristics~\citep{bandt2019small}.

The \num{1820} citations received by the seminal paper appeared in \num{684} journals indexed by the Web of Science.
These journals belong to \num{127} categories, spanning from Multidisciplinary Physics (\SI{24}{\percent} of the publications) to Zoology (only one of the of citing articles).
There are \num{17} citing articles from journals that belong to the Statistics \& Probability category.
Five of these articles appeared in \textit{Stochastic Environmental Research and Risk Assessment}, 
two in the \textit{Journal of Time Series Analysis}, 
and each of the remaining ten appeared in a different journal.
Most of these articles relate successful applications of the Bandt and Pompe methodology, except \cite{OrdinalPatternProbabilities} that obtained the sample entropy's properties under zero-mean Gaussian processes.
It is also noteworthy that, in this category of publications, \cite{DistributionsofOrderPatternsofIntervalMaps} provided a formal and more general proof of the structure of the boundary of the $H\times C$ manifold than that obtained by \cite{martin2006generalized}.
The lack of attention that the Bandt and Pompe approach has received by the Probability \& Statistics community confirms that it is a fertile research avenue waiting to explore.

In this context, the analysis of ordinal patterns coupled with the use of information theory descriptors, in addition to meeting the requirements above, has been able to detect causal information related to the unobserved variables that control the system, in addition to identifying chaotic components, visualization and characterization of different dynamic regimes, among other applications.

\section{Ordinal Patterns Representation}\label{sub:ordinalPatterns}

\cite{Bandt2002Permutation} introduced the representation of time series by ordinal patterns as a transformation resistant to noise, and invariant to nonlinear monotonic transformations.
Let ${\mathcal X} \equiv \{x_t\}_{t=1}^{T}$ be a real valued time series of length $T$, without ties. 
As stated by \cite{Bandt2002Permutation} in their seminal work:  
\begin{quote}
``If the $\{x_t\}_{t=1}^{T}$ attain infinitely many values, it is common to replace them by a symbol sequence 
$\Pi \equiv \{\pi_j\}$ with finitely many symbols, and calculate source entropy from it".
\end{quote}
Also, as stressed by these authors, 
\begin{quote}
``The corresponding symbol sequence must come 
naturally from the $\{x_t\}_{t=1}^{T}$ without former model assumptions".
\end{quote}

Let ${\mathfrak A}_{D}$ (with $D \geq 2$ and $D \in {\Bbb N}$) be the symmetric group of order $D!$ formed by all 
possible permutation of order $D$, and the symbol component vector 
${\bm \pi}^{(D)} = (\pi_1, \pi_2, \dots, \pi_D)$ so every element ${\bm \pi}^{(D)}$ is unique 
($\pi_j \neq \pi_k$ for every $j \neq k$). 
Consider for the time series ${\mathcal X} \equiv \{x_t\}_{t=1}^{T}$ its time delay embedding representation,
with embedding dimension $D \geq 2$ and time delay $\tau \geq 1$ ($\tau \in {\Bbb N}$, also called ``embedding time,'' ``time delay'', or ``delay''):
\begin{equation} 
\label{SAR:eq:time-delay}
{\mathbf X}^{(D,\tau)}_t =( x_t,x_{t+\tau},\dots,x_{t+(D-1)\tau} ) ,
\end{equation} 
for $t = 1,2,\dots,N$ with $N = T-(D-1) \tau$.
Then the vector ${\mathbf X}^{(D,\tau)}_t$ can be mapped to a symbol vector ${ \widetilde{\bm \pi}}_t^D \in {\mathfrak A}_{D}$. 
This mapping is such that preserves the desired relation between the elements 
$x_t  \in {\mathbf X}^{(D,\tau)}_t$, and all $t \in \{1,\dots,T-(D-1)\tau\}$ that share this pattern (also called ``motif'') are mapped to the same 
${\widetilde{\bm \pi}}_t^{D}$.

We define the mapping ${\mathbf X}_t^{(D,\tau)} \mapsto {\widetilde{\pi}}_t^{D}$ by ordering the observations $x_t \in {\mathbf X}_t^{(D,\tau)}$ in increasing order.
Consider the time series:
\begin{equation}
    \mathcal X = (1.8, 1.2, 3.2, 4.8, 4.2, 4.5, 2.3, 3.7, 1.2, 0.5)
    \label{eq:series}
\end{equation}
depicted in Fig.~\ref{SAR:Fig:IntroBP}.
Assume we are using patterns of length $D=5$ with unitary time lag $\tau=1$.
The code associated to $\mathbf X_{3}^{(5,1)}=(x_3,\dots,x_7)=(3.2, 4.8, 4.2, 4.5, 2.3)$, shown in black, is formed by the indexes in $\bm\pi_3^{5}=(1,2,3,4,5)$ which sort the elements of $\mathbf X_{3}^{(5,1)}$ in increasing order: $51342$.
With this, $\widetilde{\pi}_3^{5} = 51342$, and we increase the counting related to this motif in the histogram of all possible patterns of size $D=5$.

The dash-dot line in Fig.~\ref{SAR:Fig:IntroBP} illustrates $\mathbf X_{1}^{(5,2)}$, i.e. the sequence of length $D=5$ starting at $x_1$ with lag $\tau=2$.
In this case, $\mathbf X_{1}^{(5,2)}= (1.8, 3.2, 4.2, 2.3, 1.2)$, and the corresponding motif is $\widetilde{\pi}_1^{5}=51423$.

\begin{figure}
	\centering
	\includegraphics[width=.9\linewidth]{Figures/IntroBP.pdf}
	\caption{Illustration of the Bandt and Pompe coding\label{SAR:Fig:IntroBP}. Here the gray line represents the analyzed sequence $\mathcal X = (1.8, 1.2, 3.2, 4.8, 4.2, 4.5, 2.3, 3.7, 1.2, 0.5)$, the sequence illustrated by the dotted line shows the path taken when applying $\tau = 2$, and the sequence illustrated by the black line shows the elements of the pattern $\mathbf X_{1}^{(5,2)}= (1.8, 3.2, 4.2, 2.3, 1.2)$.}
\end{figure}

The classic approach to calculating the probability distribution of ordinal patterns is through the frequency histogram.
Denote $\Pi$ the sequence of symbols obtained by a given series $\mathbf{X}_t^{(D,\tau)}$.
The Bandt-Pompe probability distribution is the relative frequency of symbols in the series against the $D!$ possible patterns $\{\widetilde\pi_t^D \}_{t = 1}^{D!}$:
\begin{equation}
p(\widetilde\pi_t^D) = \frac{\#\left \{\mathbf{X}_t^{(D,\tau)} \text{ is of type } \widetilde\pi_t^D\right \}}{T- (D-1)\tau},  
\end{equation}
where  $t\in \{1, \dots, T-(D-1)\tau\}$.
These probabilities meet the conditions $p(\widetilde\pi_t^D) \ge 0$ and  $\sum_{i=1}^{D!} p(\widetilde\pi_t^D) = 1$, and are invariant before monotonic transformations of the time series values.
For example, the presence of $\alpha$ multiplicative noise in ${\mathcal X}$ does not change the results of the patterns produced.

In the literature, there are two ways to perform the mapping ${\mathbf X}_t^{(D, \tau)} \mapsto {\mathbf \pi}_t^{D}$ in the symbolization of Bandt-Pompe:~(i) Sorting the positions of the elements $x_t \in {\mathbf X}_t^{(D, \tau)}$ in chronological order (Permutation of Classification), and~(ii) Sorting the time indexes of the elements $x_t \in {\mathbf X}_t^{(D, \tau)}$ (Chronological Index Permutation).

\subsection{Rank Permutation Mapping}\label{sub:rankPermutation}

For an arbitrary $t$, the vector of real values $\mathcal{X}^{(D, \tau)}_t = (x_t, x_{t + 1}, \dots, x_{t + D - 1})$ with time delay $\tau \geq 1$ $(\tau \in \mathbb{Z})$, and embedding dimension $D \geq 2$ $(D \in \mathbb{Z})$ are mapped to the vector symbol $\bm \pi^D = (\bm R[x_t], \bm R[x_{t + 1}], \dots, \bm R[x_{t + D - 1}]) \in {\mathfrak A}_{D}$ formed by the rank of its components, defined as the following function:

\begin{equation}
    R[x_{t + n}] = \sum^{D - 1}_ {k = 0} \mathds{1} (x_{t + k} \leq x_{t + n}) \text{ for } n = 0, \dots, D - 1
\end{equation}
where $x_{t+n} \in X^{(D, \tau)}_t$, $1 \leq R(x_{t+n}) \leq D$, and $\mathds{1}$ is the indicator function: $\mathds{1}(Z) = 1$ if $Z$ is true and $0$ otherwise.  
So the maximum and minimum values of rank are $R(\min(x_{t+k})) = 1$ and $R(\max(x_{t+k})) = D$.
The complete alphabet is all the possible permutation of the ranks.
Examples of its application can be seen in~\cite{riedl2013practical} and~\cite{bandt2007order}.

For example, let us take the series~\ref{eq:series}, time delay $\tau = 1$ and embedding dimension $D = 3$, the eight vectors and their corresponding patterns obtained by rank permutation mapping are:

\begin{center}
    \scalebox{.8}{
    \begin{minipage}[b]{0.5\textwidth}
        \raggedright
        \begin{equation*}
            \bm X^{(3, 1)}_1 = (1.8, 1.2, 3.2) \mapsto (2,1,3),
        \end{equation*}
        \begin{equation*}
            \bm X^{(3, 1)}_2 = (1.2, 3.2, 4.8) \mapsto (1,2,3),
        \end{equation*}
        \begin{equation*}
            \bm X^{(3, 1)}_3 = (3.2, 4.8, 4.2) \mapsto (1,3,2),
        \end{equation*}
        \begin{equation*}
            \bm X^{(3, 1)}_4 = (4.8, 4.2, 4.5) \mapsto (3,1,2),
        \end{equation*}
    \end{minipage}
    \begin{minipage}[b]{0.5\textwidth}
        \raggedright
        \begin{equation*}
            \bm X^{(3, 1)}_5 = (4.2, 4.5, 2.3) \mapsto (2,3,1),
        \end{equation*}
        \begin{equation*}
            \bm X^{(3, 1)}_6 = (4.5, 2.3, 3.7) \mapsto (3,1,2),
        \end{equation*}
        \begin{equation*}
            \bm X^{(3, 1)}_7 = (2.3, 3.7, 1.2) \mapsto (2,3,1),
        \end{equation*}
        \begin{equation*}
            \bm X^{(3, 1)}_8 = (3.7, 1.2, 0.5) \mapsto (3,2,1),
        \end{equation*}
    \end{minipage}}
\end{center}

In Fig.~\ref{fig:Rank}, we present an illustrative drawing of this mapping for all alternatives when we have $D = 3$.
As we can see, the vertical axis indexes are fixed and the resulting pattern is obtained through the time axis labels in chronological order.

\begin{figure}
    \centering
    \includegraphics[width=\linewidth]{Figures/rank.pdf}
    \caption{\textbf{Rank permutation mapping:} The complete alphabet for $D = 3$ of the rank mapping technique obtained by permuting all possible ranks.}
    \label{fig:Rank}
\end{figure}

\subsection{Chronological Index Permutation Mapping}\label{sub:chronological}

Again, for an arbitrary $t$, the vector of real values $\mathcal{X}^{(D, \tau)}_t = (x_t, x_{t + 1}, \dots, x_{t + D - 1})$ with time delay $\tau \geq 1$ $(\tau \in \mathbb{Z})$, and embedding dimension $D \geq 2$ $(D \in \mathbb{Z})$ are mapped to the vector symbol $\bm \pi^D = (\bm i_1, \bm i_2, \dots, \bm i_D) \in {\mathfrak A}_{D}$ formed by  the time indexes ordered according to their amplitudes. 
So, the sequence must comply $x_{t+i_1} < x_{t+i_2} < \dots < x_{t+i_D}$. 
Examples of its application can be seen in~\cite{Bandt2002Permutation}, ~\cite{DistinguishingNoiseFromChaos},~\cite{parlitz2012classifying}, and~\cite{bian2012modified}.

For example, let us take the series~\ref{eq:series}, time delay $\tau = 1$ and embedding dimension $D = 3$, the eight vectors and their corresponding patterns obtained by rank permutation mapping are:

\begin{center}
    \scalebox{.8}{
    \begin{minipage}[b]{0.5\textwidth}
        \raggedright
        \begin{equation*}
            \bm X^{(3, 1)}_1 = (1.8, 1.2, 3.2) \mapsto (2,1,3),
        \end{equation*}
        \begin{equation*}
            \bm X^{(3, 1)}_2 = (1.2, 3.2, 4.8) \mapsto (1,2,3),
        \end{equation*}
        \begin{equation*}
            \bm X^{(3, 1)}_3 = (3.2, 4.8, 4.2) \mapsto (1,3,2),
        \end{equation*}
        \begin{equation*}
            \bm X^{(3, 1)}_4 = (4.8, 4.2, 4.5) \mapsto (2,3,1),
        \end{equation*}
    \end{minipage}
    \begin{minipage}[b]{0.5\textwidth}
        \raggedright
        \begin{equation*}
            \bm X^{(3, 1)}_5 = (4.2, 4.5, 2.3) \mapsto (3,1,2),
        \end{equation*}
        \begin{equation*}
            \bm X^{(3, 1)}_6 = (4.5, 2.3, 3.7) \mapsto (2,3,1),
        \end{equation*}
        \begin{equation*}
            \bm X^{(3, 1)}_7 = (2.3, 3.7, 1.2) \mapsto (3,1,2),
        \end{equation*}
        \begin{equation*}
            \bm X^{(3, 1)}_8 = (3.7, 1.2, 0.5) \mapsto (3,2,1),
        \end{equation*}
    \end{minipage}}
\end{center}

In Fig.~\ref{fig:Chronological}, we present an illustrative drawing of this mapping for all alternatives when we have $D = 3$.
In this mapping, the time indexes are fixed in chronological order.
The patterns are chosen by the vertical axis labels in the increasing direction of the amplitude.

\begin{figure}
    \centering
    \includegraphics[width=\linewidth]{Figures/Chronological.pdf}
    \caption{\textbf{Chronological index permutation mapping:} The complete alphabet for $D = 3$ of the chronological mapping technique obtained by permuting all of its indexes.}
    \label{fig:Chronological}
\end{figure}

\section{Ordinal Patterns Transition Graphs}\label{sub:OPTG}

Alternatively, given the sequence of ordinal patterns $\Pi$ one may form an oriented graph with the transitions from $\widetilde\pi_t^D$ to $\widetilde\pi_{t+1}^D$. 
The Ordinal Pattern Transition Graph ${G} = ({V}, {E})$ 
represents the transitions between two consecutive ordinal patterns over time $t$.
The vertices are the $D!$ possible patterns for an embedding dimension $D$, and the edges the transitions between them:
$V = \{v_{\widetilde\pi_t^D}\}$, and 
$E = \{(v_{\widetilde\pi_t^D}, v_{\widetilde\pi_{t+1}^D}): v_{\widetilde\pi_t^D}, v_{\widetilde\pi_{t+1}^D} \in V \}$~\citep{Borges2019Transition}.

The literature reports two approaches to compute the weight of edges.
Some authors employ unweighted edges~\citep{McCullough2015lagged, Kulp2016ordinal}, which represent only the existence of transitions, while others apply the frequency of transitions~\citep{Sorrentino2015periodic, Zhang2017ConstructingOP}.
The weights $\mathbbm{W} = \{w_{v_{\widetilde{\pi}^D_i}, v_{\widetilde\pi^D_j}}: v_{\widetilde\pi^D_i}, v_{\widetilde\pi^D_j} \in V \}$ assigned to each edge describe the chance of transitions between the patterns $(v_{\widetilde\pi^D_i}, v_{\widetilde\pi^D_j})$
The weights are calculated as the relative frequency of each transition, i.e.:
\begin{equation}
w_{v_{\widetilde\pi^D_i}, v_{\widetilde\pi^D_j}} = \frac{|\Pi_{\widetilde\pi^D_i,\widetilde\pi^D_j}|}{T-(D-1)\tau-1},
\end{equation}
where $|\Pi_{\widetilde\pi^D_i,\widetilde\pi^D_j}|$ is the number of transitions from pattern $\widetilde\pi^D_i$ to pattern $\widetilde\pi^D_j$, $\sum_{v_{\widetilde\pi^D_i}, v_{\widetilde\pi^D_j}}w_{v_{\widetilde\pi^D_i}, v_{\widetilde\pi^D_j}} = 1$,
and the denominator is the number of transitions between sequential patterns in the series of motifs of length $T-(D-1)\tau$.

When comparing the transition graph with other classical time series representations in graphs, we can highlight some advantageous properties~\citep{Borges2019Transition}:
\begin{enumerate}[label=(\roman*)]
    \item \textbf{Speed and simplicity}: the construction of the graph depends only on two parameters, the time delay $\tau$ and the embedding dimension $D$.
    On the other hand, the time series symbolization process depends on the length $N$ of the analyzed series and the dimension of incorporation $D$, having a complexity $\mathcal{O} (n \cdot D^2)$ when we consider a simple algorithm for sorting elements with $\mathcal{O}(D^2)$ complexity.
    \item \textbf{Scalability}: Regardless of the number of elements in the time series, the number of vertices of the transition graph will always be limited by $D!$, Where $3 \leq D \leq 7$.
    \item \textbf{Robustness}: Since the Bandt-Pompe symbolization produces robust ordinal patterns to the presence of monotonous non-linear noises and transformations, we can conclude that the transformation proposed by the transition graphs also guarantees such specificity.
    \item \textbf{Probability of self-transition:} The self-transition of the graph represent the proportion of ordinal patterns sequentially repeated and are defined as
    \begin{equation}
        p_{st} = p_{(\widetilde \pi^D_i, \widetilde \pi^D_i)} = \sum_{i \in \{1, \cdot, D!\}} w_{v_{\widetilde\pi^D_i}, v_{\widetilde\pi^D_i}}.
    \end{equation}
    Thus, the presence of auto-loops, when evaluated according to the $\tau$ incorporation delay, proved to be an important characteristic of the underlying dynamics of time series, as they are directly associated with the temporal correlation of the elements.
\end{enumerate}

\section{Information-Theoretic Descriptors}\label{sub:InformationTheory}

Let $\mathbbm{P} = \{p_{\widetilde\pi^D_1}, p_{\widetilde\pi^D_2}, \dots, p_{\widetilde\pi^D_{D!}} \} = \{p_1,\dots,p_{D!}\}$ be the probability function obtained from the \mbox{1-D} signal $\mathcal{X}$ by Bandt-Pompe symbolization.
The last step of the characterization process consists of calculating the Information Theory descriptors: Shannon Entropy and Statistical Complexity.
Through these features we were able to obtain the point in the $H \times C$ plane.

\subsection{Permutation Entropy}

Entropy measures the disorder or unpredictability of a system characterized by a probability measure $\mathbbm{P}$ and is measured by:
\begin{equation}
H(\mathbbm{P}) = -\sum_{i = 1}^{D!} p_{\widetilde\pi^D_i} \log p_{\widetilde\pi^D_i} .
\label{eq:Entropia}
\end{equation}

Its minimum value occurs when $H (\mathbbm{P}) = H_{min} = 0$, in this particular case we will have maximum knowledge about the system and we will be able to predict with absolute certainty what will be the next ordinal pattern generated by the data, indicating that the time series is deterministic.
On the other hand, when the behavior of the system is described by a uniform distribution, that is, when all the possibilities have the same probability of occurrence and its probability is determined by $\mathbb{P} = \{1/D!, \dots, 1/D! \}$, we will have minimal knowledge of the analyzed data, obtaining $H(\mathbbm{P}) = H_{max} = \log D!$.
In that case, the time series would be determined as a completely random system~\citep{Bandt2002Permutation}.

However, in the literature it is opted to use the normalized Shannon entropy defined by~\cite{MARTIN2006439}, given by:
\begin{equation}
    H_S(\mathbb{P}) = \frac{H(\mathbb{P})}{H_{max}} = - \frac{1}{\log D!}\sum_{i = 1}^{D!} p_{\widetilde\pi^D_i} \log p_{\widetilde\pi^D_i},
\end{equation}
where $0 \leq H_S(\mathbb{P}) \leq 1$.

\subsection{Statistical Complexity}

The entropy's ability to capture system properties is limited, so it is necessary to use it in conjunction with other descriptors to obtain a complete analysis.
Other interesting measures are the distances between $\mathbbm{P}$ and a probability measure that describes a non-informative process, typically the uniform distribution.

The Jensen-Shannon divergence to the uniform distribution $\mathbbm{U} = (\frac{1}{D!}, \dots, \frac{1}{D!})$ is a measure of how similar the underlying dynamics is to a non-informative process.
It is calculated as:
\begin{equation}
    JS(\mathbbm{P}, \mathbbm{U}) = \sum_{i = 1}^{D!} \Big(p_{\widetilde\pi^D_i} \log\frac{p_{\widetilde\pi^D_i}}{u_{\widetilde\pi^D_i}} +
    u_{\widetilde\pi^D_i} \log\frac{u_{\widetilde\pi^D_i}}{p_{\widetilde\pi^D_i}}
    \Big).
\end{equation}

Conversely to entropy, the statistical complexity seeks to find interaction and dependence structures among the elements of a given series, being an extremely important factor in the study of dynamic systems.
The Statistical Complexity is defined as~\cite{Lamberti2004Entropic}:
\begin{equation}
    C_{JS}(\mathbbm{P}, \mathbbm{U}) = H_S(\mathbbm{P}) \cdot Q_{JS}(\mathbbm{P}, \mathbbm{U}).
\end{equation}

The ``disequilibrium'' $Q_{JS}$ is calculated by:
\begin{equation}
    Q_{JS}(\mathbbm{P}, \mathbbm{U}) = Q_0 \cdot JS(\mathbbm{P}, \mathbbm{U})
\end{equation}
\begin{equation}
    = Q_0 \cdot \left\{ H \left[ \frac{\mathbbm{P} + \mathbbm{U}}{2} \right] - \frac{H[\mathbbm{P} + \mathbbm{U}]}{2} \right\}
\end{equation}
The normalization constant is equal to the inverse of the maximum value of $JS$~\citep{rosso2007extracting}, so:
\begin{equation}
    Q_0 = -2 \left\{ \left( \frac{D! + 1}{D!} \right) \ln{(D! + 1)} - 2\ln{(2D!)} + \ln{(D!)} \right\}^{-1},
\end{equation}
where $0 \leq Q_0 \leq 1$.

This descriptor quantifies, in addition to randomness, the presence of correlational structures between patterns, reflecting the architecture of the systems, being different from zero if there are more likely states than the others.
In this way, different degrees of structure can be quantified, reflecting properties revealed by the probability distribution of the underlying process.

\subsection{Causality Complexity-Entropy Plane}

The set of all pairs $(H(\mathbbm{P}), C(\mathbbm{P}, \mathbbm{U}))$ for any time series described by patterns of length $D$ lies in a compact subset of $\mathbbm R^2$: the Complexity-Entropy plane (or Entropy-Complexity plane).
Through this tool it is possible to discover the nature of the series just by checking its region of location on the plane, its associated values help to determine whether it corresponds to a chaotic (or other deterministic dynamics) or stochastic sequence.

\citet{SomeFeaturesoftheLMCStatisticalComplexity} proved that, for a fixed value of entropy, there are two extreme values of complexity.
\citet{GeneralizedStatisticalComplexityMeasuresGeometricalAnalyticalProperties}, using geometrical arguments on the space of configurations, found expressions for such boundaries.
The lower boundary $C_{\min}$ is continuous, while the upper $C_{\max}$ is defined by $D!-1$ pieces.

Fig.~\ref{fig:Boundaries} shows the boundaries of the $H\times C$ plane for the embedding dimensions $D=3$ (red) $D=4$ (green), and $D=5$ (blue).
The jagged structure of $C_{\max}$ increases the difficulty of finding distributions for the points in the $H\times C$ plane.

\begin{figure}[hbt]
    \centering
    \includegraphics[width=\linewidth]{Figures/BoundariesPlot}
    \caption{Boundaries of the $H\times C$ plane for $D=3,4,5$.}\label{fig:Boundaries}
\end{figure}

We illustrate the use of the Complexity-Entropy ($H\times C$) with the following time series:
\begin{itemize}
    \item \textbf{$f^{-k}$ noises}. Correspond to colored noise sequences generated synthetically using the power spectrum $f^{-k}$~\citep{DistinguishingNoiseFromChaos}.
    Here, we opt to use the following configurations of the spectrum: white ($k=0$), $k=-1/2$, pink ($k=1$), $k=3/2$, red ($k=2$), $=5/2$, and $k=3$;
    \item \textbf{Logistic maps}. The classical chaotic map defined by the formula $x_t = r x_{t-1} (1 - x_{t-1})$~\citep{peitgen2006chaos}.
    In this work, we use this chaotic logistic series with the parameters $r = 3.6$ and $4$;
    \item \textbf{Deterministic series}. Monotonic increasing ($\log(x_t+0.1)$, $x_t=\{1,2,\dots,10^4$) and periodic ($\sin(2x_t)\cos(2x_t)$, with $0\leq x_t\leq 2\pi$ over ten thousand equally spaced points).
\end{itemize}
In all cases, we used $D=6$ and $\tau=1$.
Fig.~\ref{fig:Histograms} shows nine of the histograms produced by these series using the Mersenne-Twister pseudorandom number generator;
we omitted those corresponding to the deterministic series, as they produce one and two nonzero bins.
As we can see, as we add more correlation structures, $f^{-k}$ noises have a less uniform histogram, making it more evident that some ordinal patterns stand out and are more frequent than others.
Therefore, we have an indication that there is a structure of dependence between such elements, making them more deterministic.

\begin{figure}
    \includegraphics[width=\linewidth]{Figures/h.pdf}
	\caption{Patterns histograms of selected time series with $D=6$, $\tau=1$ and $T = \num[scientific-notation=true]{e4}$.}
	\label{fig:Histograms}
\end{figure}

Fig.~\ref{fig:AllSystems} shows the $H\times C$ plane with the bounds for $D=6$, the time series, and the points they were mapped onto.
The points due to $f^{-k}$ noises appear joined by dotted segments.
It is noticeable that deterministic patterns have more complexity than random ones.
Also, points related to $f^{-k}$ noises tend to clutter for $k<1$, having the highest entropy values.

\begin{figure}
    \centering
    \includegraphics[width=\linewidth]{Figures/AllSystems.pdf}
    \caption{Eleven systems and their points in the $H\times C$ plane when we apply $D = 6$, $\tau = 1$ and $T = \num[scientific-notation=true]{e4}$.}
    \label{fig:AllSystems}
\end{figure}

	


\chapter{Weighted Amplitude Transition Graph}\label{chapter:WATG}

Although the use of Bandt-Pompe symbolization has several advantages over other feature extraction algorithms, it has two major gaps in its seminal definition: 
\begin{enumerate}[label=(\roman*)]
    \item  The ordinal ambiguity present when we have equal values in the same sub-sequence~\citep{Traversaro2018DDMI,Traversaro2017Empirical}, and
    \item The lack of information related to differences in sample amplitude~\citep{azami2016amplitude, Fadlallah2013Weightedpermutation, cuesta2019permutation},that is, the mean value of the amplitudes and the differences between neighboring samples are not considered by the original methodology.
\end{enumerate}
Figure~\ref{fig:possible-patterns} shows two examples of how sub-sequences with different structural characteristics can be mapped to the same ordinal pattern, thus decreasing the characterization power of the approach.

By adding amplitude information, we were able to attribute less complexity in sequences with greater regularity and locate abrupt changes along the signal, reducing the impact of a possible noise degradation in relation to the final value of the descriptors.
In this context, several modifications were proposed in the calculation of the Bandt-Pompe distribution with the aim of neutralizing the limitations discussed and maintaining most of the properties presented in the technique initially.

To analyze the different methods present in the literature, we will apply the following sequence as a numerical example:
\begin{equation}
    \mathcal{X} = \{-3.7, -3.5, 2, 1.3, 0.8, -2.3, 1.8, 1.7, 1.3, 2.6, 1.7, 0.9, 0, -0.4, -0.5, 7\}
	\label{eq:example}
\end{equation}
The Chronological Index Permutation was used to generate the ordinal patterns, where for $D = 3$ and $\tau = 1$, we will have $M = 15$ vectors $\mathbf{X}^{(D = 3)}$ .
Its probability distribution obtained by the traditional Bandt-Pompe symbolization method is
$\mathbb{P} = \{0, 0.2, 0.2, 0.0666667, 0.0666667, 0.4666667\}$.

\begin{figure}
	\centering
	\includegraphics[width=\textwidth]{Figures/possible-patterns.pdf}
	\caption{Two examples of possible $D$-dimensional vectors corresponding to the same ordinal patterns when using $D = 3$.}
	\label{fig:possible-patterns}
\end{figure}

\section{Weighted Ordinal Patterns Methods in Literature}\label{Methods}

There are basically two strategies for incorporating amplitude information into ordinal structures.
The first assumes that the greater the variation in amplitude in a given ordinal pattern, the greater its weight in calculating the probability distribution.
Thus, correction factors are proposed to carry out such weighting.
The second proposes an extension in the alphabet of symbols, mapping vectors with different amplitudes to different ordinal patterns.
Thus, the total set of patterns has more than $D!$ Possible components and vary according to the mapping method considered (rank permutation or chronological index permutation).

\subsection{Weighted Permutation Entropy}\label{WPE}

The Weighted Permutation Entropy (WPE) was proposed by ~\cite{Fadlallah2013Weightedpermutation}. 
The proposed idea is to weigh the relative frequencies of the patterns taking into account the variability of the elements in relation to the average value of the analyzed segment.
Denote $\overline{X}_t^{(D, \tau)}$ the arithmetic mean:
\begin{equation}
\overline{X}_t^{(D, \tau)} = \frac{1}{D} \sum_{k = 1}^{D} x_{t + (k - 1)}.
\end{equation}
The weight $w_{t}$ is the sample variance of each vector $X_t^{(D, \tau)}$:
\begin{equation}
w_{t} = \frac{1}{D} \sum_{k = 1}^{D}\big[x_{t + (k - 1)} - \overline{X}_t^{(D, \tau)}\big]^2 .
\end{equation}
Then, the probability distribution is given from the weighted relative frequencies:
\begin{equation}
p(\widetilde \pi_t^D) = \frac{\sum_{i : \{\mathbf{X}^{(D,\tau)}_i \mapsto \widetilde\pi^D_t\}} w_{i}}{\sum_{i = 1}^{T-(D-1)\tau} w_{i}}.
\end{equation}

Using the time series~\ref{eq:example} as an example, the weight shown by WPE for each subsequence (the variance of each vector) would be as follows:
$w_1 = 23.389$, 
$w_2 = 18.389$,
$w_3 = 5.976$,
$w_4 = 0.242$,
$w_5 = 2.536$,
$w_6 = 3.047$,
$w_7 = 3.647$,
$w_8 = 0.047$,
$w_9 = 0.296$,
$w_{10} = 0.296$,
$w_{11} = 0.482$,
$w_{12} = 0.482$,
$w_{13} = 0.296$,
$w_{14} = 0.047$,
$w_{15} = 12.336$.
Being described in the $H \times C$ plane by the following point
$(\bm h, \bm c) = (0.965, 0.0322)$.

\subsection{Fine-Grained Permutation Entropy}\label{FGPE}

The Fine-Grained Permutation Entropy (FGPE) was introduced in Ref.~\cite{xiao2009fine}.
Let $\beta_t$ be the difference series:
\begin{equation}
\beta_t = \big\{|x_{t+1} - x_t|, \dots, |x_{t+(D-1)} - x_{t+(D-2)}|\big\}.
\end{equation}
The weight $w_t$ quantifies such differences:
\begin{equation}
w_t = \left \lfloor \frac{\max\{\beta_t\}}{\alpha s(\beta_t) } \right \rfloor,
\end{equation}
where $s$ is the sample standard deviation,
$\alpha$ is a user-defined parameter, 
and $\lfloor\cdot \rfloor$ is the floor function.
Then, $w_t$ is added as a symbol at the end of the corresponding pattern to quantify the difference between neighboring elements in the vector, leading to an update of $\Pi$:
\begin{equation}
\pi{'}_t^D = \{ \widetilde \pi_t^D \cup w_t\}.
\end{equation} 
Finally, the probability distribution is calculated as:
\begin{equation}
p(\pi{'}_t^D) = \frac{\#\left \{\mathbf{X}_t^{(D,\tau)} \text{ is of type } \pi{'}_t^D\right \}}{T- (D-1)\tau}.
\end{equation}
Using the time series~\ref{eq:example} as an example, when applying FGPE we will have the formation of the following symbols:

\begin{center}
    \scalebox{.75}{
    \begin{minipage}[b]{0.25\textwidth}
        \raggedright
        $$\bm{X}^{(3)}_1 	\mapsto (2,0,1,59),$$
        $$\bm{X}^{(3)}_2 	\mapsto (1,2,0,5),$$
        $$\bm{X}^{(3)}_3 	\mapsto (0,2,1,3),$$
    \end{minipage}
    \begin{minipage}[b]{0.25\textwidth}
        \raggedright
        $$\bm{X}^{(3)}_4 	\mapsto (2,1,0,9),$$
        $$\bm{X}^{(3)}_5 	\mapsto (2,1,0,3),$$
        $$\bm{X}^{(3)}_6 	\mapsto (1,0,2,11),$$
    \end{minipage}
    \begin{minipage}[b]{0.25\textwidth}
        \raggedright
        $$\bm{X}^{(3)}_7 	\mapsto (0,2,1,0),$$
        $$\bm{X}^{(3)}_8 	\mapsto (2,1,0,3),$$
        $$\bm{X}^{(3)}_9 	\mapsto (1,0,2,4),$$
    \end{minipage}
    \begin{minipage}[b]{0.25\textwidth}
        \raggedright
        $$\bm{X}^{(3)}_{10} \mapsto (0,2,1,9),$$
        $$\bm{X}^{(3)}_{11} \mapsto (2,1,0,25),$$
        $$\bm{X}^{(3)}_{12} \mapsto (2,1,0,25),$$
    \end{minipage}
    \begin{minipage}[b]{0.25\textwidth}
        \raggedright
        $$\bm{X}^{(3)}_{13} \mapsto (2,1,0,5),$$
        $$\bm{X}^{(3)}_{14} \mapsto (2,1,0,3),$$
        $$\bm{X}^{(3)}_{15} \mapsto (1,0,2,2)$$
    \end{minipage}}
\end{center}

\subsection{Amplitude-Aware Permutation Entropy}\label{AAPE}

The Amplitude-Aware Permutation Entropy (AAPE) was proposed by~\cite{azami2016amplitude} and presents a solution to the absence of amplitude information and also quantifies the possible effect of repeated values during symbolization.
Here we will report only its simplified version presented in~\cite{cuesta2018classification}.
It consists of weighting the amplitude of ordinal patterns by both the mean and the differences of the elements.
For this, only an additional parameter $ A \in [0, 1]$ is required:
\begin{equation}
w_t = \frac{A  |x_t|}{D} + \sum_{k = 1}^{D - 1} \left(\frac{A  |x_{t + k}|}{D} + \frac{(1 - A)  |x_{t+k} - x_{t+k-1}|}{D - 1}\right).
\end{equation}
The AAPE consists of one of the methods present in the first approach mentioned, that is, it will also apply a correction factor to the ordinal patterns obtained by the analyzed sequence.
In this way, the probability histogram, instead of being updated each time a match is found, will be updated using the weighted relative frequency:
\begin{equation}
p(\widetilde \pi_t^D) = \frac{\sum_{i : \{\mathbf{X}^{(D,\tau)}_i \mapsto \widetilde\pi^D_t\}} w_{i}}{\sum_{i = 1}^{T-(D-1)\tau} w_{i}}.
\end{equation}

Using the time series~\ref{eq:example} as an example, given the ordinal patterns obtained, the AAPE algorithm acquires the following weights for each subsequence:
$w_1 = 2.683$, 
$w_2 = 2.016$,
$w_3 = 1.483$,
$w_4 = 2.616$,
$w_5 = 1.358$,
$w_6 = 3.216$,
$w_7 = 6.083$,
$w_8 = 7.375$,
$w_9 = 0.983$,
$w_{10} = 1.633$,
$w_{11} = 0.925$,
$w_{12} = 1.291$,
$w_{13} = 0.858$,
$w_{14} = 0.541$,
$w_{15} = 0.275$.

\section{Proposed Method}\label{WATG}

Recent approaches to extracting knowledge in time series have been developed using information from transition graphs of ordinal patterns.
However, such tooling comes up against the same problem as the traditional symbolization of Bandt-Pompe, the lack of information of amplitude between different sequences.
In order to be the first variation of transition graphs to incorporate this type of information, we propose in this present work the Weighted Amplitude Transition Graph (WATG).

Since we are interested in a comparable metric for different data sets, the first step is to normalize each time series $\mathcal{X}$ to the interval $[0, 1]$:
\begin{equation}
\frac{x_i - x_{\min}}{x_{\max} - x_{\min}} \longmapsto x_i,
\label{eq:scaling}
\end{equation}
where $x_{\min}$ and $x_{\max}$ are, respectively, the minimum and maximum values of the series.
This transformation is relatively stable before contamination, e.g., if instead of $x_{\max}$ we observe $k x_{\max}$ with $k\geq 1$, the relative values are not altered. Nevertheless, other more resistant transformations as, for instance, $z$ scores, might be considered.

Assuming that a variation in amplitude in a given ordinal pattern should directly impact the resulting probability distribution.
Thus, we assign for each vector $\mathbf{X}^{(D, \tau)}_t$ an associated weight $\beta_t$ that measures the biggest difference between its elements:
\begin{equation}
\beta_t = \max\{|x_i - x_j|\},
\end{equation}
where $x_i, x_j \in \mathbf{X}^{(D, \tau)}_t$.

Thus, each edge will have its value no longer associated with the number of transitions made between ordinal patterns, but with the difference in amplitude observed during the respective transition:
\begin{equation}
w_{v_{\widetilde \pi^D_i}, v_{\widetilde \pi^D_j}} =  \sum_{i : \{\mathbf{X}^{(D,\tau)}_t \mapsto \widetilde\pi^D_i\}} \sum_{j : \{\mathbf{X}^{(D,\tau)}_t \mapsto \widetilde\pi^D_j\}} |\beta_i - \beta_j| .
\end{equation}
Thus, the probability distribution taken from the weighted amplitude transition graph is given as follows:	
\begin{align}
&\left\{\begin{array}{l}
\lambda_{v_{\widetilde\pi^D_i}, v_{\widetilde\pi^D_j}} = 1, \text{ if } (v_{\widetilde\pi^D_i}, v_{\widetilde\pi^D_j}) \in {E}, \\
\lambda_{v_{\widetilde\pi^D_i}, v_{\widetilde\pi^D_j}} = 0, \text{ otherwise}.
\end{array}\right., \text{ and} \\
&p(\widetilde\pi^D_i, \widetilde\pi^D_j) = \frac{\lambda_{v_{\widetilde\pi^D_i}, v_{\widetilde\pi^D_j}} \cdot w_{v_{\widetilde\pi^D_i}, v_{\widetilde\pi^D_j}}}{\sum_{v_{\widetilde\pi^D_a}, v_{\widetilde\pi^D_b}} w_{v_{\widetilde\pi^D_a}, v_{\widetilde\pi^D_b}}}.
\end{align}

Note that  $p(\widetilde\pi^D_i, \widetilde\pi^D_j) \ge 0$ and $\sum_{\widetilde\pi^D_i, \widetilde\pi^D_j} p(\widetilde\pi^D_i, \widetilde\pi^D_j) = 1$, so $p$ is a probability function.

We were able to observe the impact of adding amplitude information to transition graphs on two different occasions:
\begin{itemize}
    \item If the signal \mbox{1-D} shows a low variation in amplitude and peak intensity between them, the transitions of the ordinal patterns that represent the latter have greater weights.
    This contributes to the probability distribution becoming less uniform, since it will be more concentrated on these edges.
    Another characteristic observed is a drop in the entropy values of these sequences when compared to the traditional method.
    \item If the signal \mbox{1-D} shows a uniform amplitude variation, the weights assigned to the edges will be well distributed, thus also giving rise to a uniform probability distribution and consequently a greater associated entropy.
\end{itemize}

The formation schematic of our weighted graph proposal is illustrated in Figure~\ref{fig:diagramWATG}, where we show how we perform the calculation of the weights of ordinal patterns.

\begin{figure}
	\includegraphics[width=\columnwidth]{Figures/watg-diagram.pdf}
	\caption{Schematic of the composition of the weighted amplitude graph (WATG).}
	\label{fig:diagramWATG}
\end{figure} 
\chapter{Case study: Classification of SAR Textures}\label{chapter:SARclassification}

We present a newly proposed technique for texture analysis and classification based on the Bandt-Pompe symbolization for SAR data.
It consists of
(i)~linearize a \mbox{2-D} patch of the image using the Hilbert-Peano curve,
(ii)~build the Ordinal Pattern Transition Graph;
(iii)~obtain a probability distribution function derived from this graph;
(iv)~compute Information Theory descriptors (Permutation Entropy and Statistical Complexity) from this distribution and use them as features to feed a classifier.
The ordinal pattern graph we propose considers that the edges' weight is related to the absolute difference of observations, which encodes the information about the data amplitude. 
This modification considers the unfavorable signal-to-noise ratio of SAR images and leads to the characterization of several types of textures.
Experiments with data from Munich urban areas, Guatemala forest regions, and Cape Canaveral ocean samples show the effectiveness of our technique in homogeneous areas, achieving satisfactory separability levels.
The two descriptors chosen in this work are easy and quick to calculate and are used as input for a $k$-nearest neighbor classifier.
Experiments show that this technique presents results similar to state-of-the-art techniques that employ a much larger number of features and, consequently, impose a higher computational cost.

\section{Introduction}\label{SAR:Intro}

Texture is an elusive trait.
When dealing with remotely sensed images, the texture of different patches carries relevant information that is hard to quantify and transform into useful and parsimonious features.
This may be since textures, in this context, is a synesthesia phenomenon that triggers tactile responses from visual inputs.
This work presents a new way of extracting features from textures, both natural and resulting from anthropic processes, in SAR (Synthetic Aperture Radar) imagery.

SAR systems are a vital source of data because they provide high-resolution images in almost all weather and day-night conditions.
They provide basilar information, complementary to that offered by sensors that operate in other regions of the electromagnetic spectrum, for a variety of Earth Observation applications.	
Although they present rich information, such data have challenging characteristics.
Most notably, they do not follow the usual Gaussian additive model, and the signal-to-noise ratio is usually low.

Yue et al.~\cite{Yue2020Gaussian} provide a comprehensive account of how the physical properties of the target are translated into first- and second-order statistical properties of SAR intensity data.
There is general agreement that non-deterministic textures are encoded in the second-order features, i.e., in the spatial correlation structure.
Therefore is frequent the use of covariance matrix and other measures that assume that a linear dependence, namely the Pearson correlation coefficient, suffices to characterize natural textures.
However, in SAR imagery, texture is often visible only over large areas, and the multiplicative and non-Gaussian nature of speckle antagonizes with the additive assumption that underlies classical approaches, making complex the process of characterizing such data.

Surface classification and land use are among the most critical applications of the Synthetic Aperture Radar (SAR) image~\cite{Pottier2004Unsupervised}.
In recent years, handcrafted features and representation learning (supervised and unsupervised) algorithms have been proposed~\cite{han2020unsupervised, huang2020classification, xie2020polsar}.
Algorithms of the unsupervised generative adversarial network (GAN) have revolutionized the classification of SAR images, improving performance in small sample problems, and helping the interpretability of such data~\cite{liu2019task}.
Among the supervised algorithms, support vector machine (SVM)~\cite{sukawattanavijit2017ga}, random forest (RF)~\cite{mcnairn2014early}, and neural network (NN)~\cite{lin2017deep} have been frequently used in remote sensing.
The Principle Component Analysis (PCA)~\cite{ressel2015neural}, autoencoder~\cite{wang2019classification} and the Boltzmann machine~\cite{qin2017object} can to extract non-local resources and classify non-labeled PolSAR pixels using an unsupervised approach.
However, methods such as graph-based semi-supervised deep learning algorithms~\cite{bi2018graph} can improve classification accuracy in problems with few labeled samples.

Handcrafted features in SAR textures can be studied following two complementary approaches, namely analyzing the marginal properties of the data (first-order statistics), and observing their spatial structure~\cite{Yue2020Gaussian, numbisi2018multi}.
In this work, we focus on the second approach, which shows relevant results using techniques from the image processing literature, such as co-occurrence matrices and Haralick's descriptors~\cite{yu2019detection}.
Through the gray-level co-occurrence matrices (GLCM), we can extract features that reflect statistical relationships of the pixel intensity values.
On the other hand, Haralick's descriptors can capture information on intensity and amplitude based on global statistics of SAR images.
Radford et al.~\cite{radford2018geological} used textural information derived from GLCM, along with Random Forests, for geological mapping of remote and inaccessible localities; the authors obtained a classification accuracy of $\approx\SI{90}{\percent}$, even when using limited training data ($\approx\SI{0.15}{\percent}$ of the total data). 	
Hagensieker and Waske~\cite{hagensieker2018evaluation} evaluated the synergistic contribution of multi-temporal L-, C-, and X-band data to tropical land cover mapping, comparing classification outcomes of ALOS-2~\cite{kankaku2013alos}, RADARSAT-2~\cite{morena2004introduction}, and TerraSAR-X~\cite{breit2009terrasar} datasets for a study site in the Brazilian Amazon using a wrapper approach. 
The wrapper utilizes the gray-level co-occurrence matrix texture information and a  Random Forest classifier to estimate scene importance. 	
Storie~\cite{storie2018urban} proposed an open-source workflow for detecting and delineating the urban-rural boundary using Sentinel-1A SAR data.
The author used a combination of GLCM information and a k-means classifier to produce a three-category map that distinguishes urban from rural areas. 
In higher resolution image classification activities, it is necessary to obtain more granular information from the data by extracting local characteristics such as scale and orientation.
In this scenario, techniques such as Fourier power spectrum~\cite{Florindo2012Fractal}, random fields~\cite{zhu2016antarctic}, Gabor filter~\cite{dumitru2014information} and wavelet transform~\cite{akbarizadeh2012new} are usually applied.

In our approach, we opt to analyze the \mbox{1-D} signals resulting from the linearization of the image samples, using non-parametric time series analysis techniques.
With this approach, we reduce the dimensionality of the data while preserving the spatial correlation structure.
Observations are then transformed into ordinal patterns with the Bandt-Pompe symbolization.
We use Information Theory descriptors to analyze the distributions these patterns induce, both directly and by building transition graphs among subsequent patterns.
Those descriptors are the Entropy and the Statistical Complexity, which are easy to obtain and are interpretable.
They reveal important features of the underlying process.

The following question guides us:
\begin{quote}
What is the best representation of a texture patch that allows extracting expressive Information Theory descriptors to characterize textures in the presence of speckle?
\end{quote}

Thus, we use the Weighted Amplitude Transition Graph (WATG), incorporating the absolute difference among observations as weights of the edges between nodes transitions.
Such weights take part in the computation of the probabilities and, thus, influence both Entropy and Statistical Complexity.

This work's main contribution is the proposal of a new representation of SAR textures, which allows a low-dimensional characterization useful for, among other applications, their classification.
We compare its performance with the classical histograms of Bandt-Pompe ordinal patterns and the regular transition graph.
Since the proposed approach has a low computational cost, the results obtained suggest that this technique has good potential in other applications, such as texture segmentation tools of SAR images.

\section{Linearization of image patches}\label{linearization}

We perform a data dimensionality reduction by turning the \mbox{2-D} patch into a \mbox{1-D} signal.
This could be accomplished by reading the data by lines, columns, or any transformation of \mbox{2-D} indexes into a sequence of integers.
In this work, we chose to use the Hilbert-Peano~\cite{Lee1994Texture} curve, due to its low computational cost and its ability to preserve relevant properties of pixel spatial correlation.

Nguyen et al.~\cite{nguyen1982space} firstly employed Space-filling curves, to map texture into a one-dimensional signal.
Carincotte et al.~\cite{Carincotte2006changeDetection} used the Hilbert-Peano curve in the problem of change detection in pairs of SAR images.
The authors noted that this transformation exploits the spatial locality and that its pseudo-randomness
of direction changes work well for a large family of images, especially
natural ones.

Assuming an image patch is supported by an $M \times N$ grid, we have the following definition.

\newtheorem{mydef}{Definition}
\begin{mydef}
	An image scan is a bijective function $f \colon \mathbbm{N} \times \mathbbm{N} \to \mathbbm{N}$ in the ordered pair set $ \{(i, j): 1 \leq i \leq M , 1 \leq j \leq N\}$, which denotes the points in the domain, for the closed range of integers $\{1, \dots, M  N\}$.
	A scan rule is $\{f^{-1}(1), \dots, f^{-1}(M  N)\}$.
	\label{def:CurveFilling}
\end{mydef}
This Definition imposes that each pixel is visited only once and that all pixels are visited.

Space-filling curves, such as raster-1, raster-2, and Hilbert-Peano scanning techniques, stipulate a proper function $f$.
Hilbert-Peano curves scan an array of pixels of dimension $2^k \times 2^k$, $k \in \mathbbm{N}$, never keeping the same direction for more than three consecutive points, as shown in Fig.~\ref{fig:Hilbert}.
Using the Hilbert-Peano curve, we reduce the data dimensionality while maintaining the patch's spatial dependence information.
In this work, we use Hilbert-Peano patches of size $128 \times 128$.

\begin{figure}
	\centering
	\tikz[scale=6.2] \hilbert((0mm,0mm),3);
	\hspace{0.3cm}
	\tikz[scale=2.95] \hilbert((0mm,0mm),4);
	\hspace{0.3cm}
	\tikz[scale=1.45] \hilbert((0mm,0mm),5);	
	\caption{Hilbert-Peano curves in areas of: (a)~$8 \times 8$, (b)~$16 \times 16$, and (c)~$32 \times 32$ pixels. }\label{fig:Hilbert}
\end{figure} 

Figs.~\ref{Fig:figGuatemala}, \ref{Fig:figCape1}, \ref{Fig:figCape2}, \ref{Fig:figMunich}, and~\ref{Fig:figPasture} show five image patches with different textures.
Figs.~\ref{Fig:TSGuatemala}, \ref{Fig:TSCape1}, \ref{Fig:TSCape2}, \ref{Fig:TSMunich}, and~\ref{Fig:TSPasture} present their 1-D representation as signals. 
	
\section{Methodology}\label{methodology}

Algorithm~\ref{alg:watg} outlines our methodology.
Line~\ref{Line:Peano} transforms the input texture patch $P$ in a \mbox{1-D} signal with a Hilbert-Peano sequence.
%
With this, the spatial information is encoded into a one-dimensional signal.
%
Line~\ref{Line:Probs} computes 
the probability distribution of the weighted transition graph induced by the \mbox{1-D} signal.
The WATG function is detailed in Lines~\ref{Line:WATGBegin}--\ref{Line:WATGEnd}.
%
Lines~\ref{Line:Shannon} and~\ref{Line:StatisticalComplexity} compute the two descriptors of the patch.

The \texttt{WATG} function consists of three steps: 
(i)~each sub-sequence of size (dimension) $D$ of observations at delay $\tau$ is transformed into an ordinal pattern using the Bandt-Pompe symbolization (function \texttt{BPSymbolization}, Line~\ref{Line:BP}); 
(ii)~function \texttt{transitions} (Line~\ref{Line:Transitions}) calculates the sequence of alternations of the ordinal patterns; and 
(iii)~function \texttt{weigthGraph} (Line~\ref{Line:WG}) generates the incidence matrix of the graph using the as weights the amplitude differences between the time series elements.
Finally, the probability distribution is obtained by turning the transition matrix into a vector (Line~\ref{Line:WATGEnd}).
These steps are also depicted in Fig.~\ref{fig:WATG}.

\begin{algorithm}
	\caption{$H \times C$ point from a patch using WATG}
	\label{alg:watg}                                
	\textbf{Input:} Patch of texture $P$, dimension $D$ and time delay \textbf{$\tau$}\\
	\textbf{Output:} $H \times C$ feature
	\begin{algorithmic}[1]
		\State signal.\mbox{1-D} $\gets$ \texttt{hilbertCurve}($P$) \label{Line:Peano}
		\State Probs $\gets$ \texttt{WATG}(signal.\mbox{1-D}, D, $\tau$) \label{Line:Probs}
		\State H $\gets$ \texttt{ShannonEntropy}(Probs) \label{Line:Shannon}
		\State C $\gets$ \texttt{StatisticalComplexity}(Probs) \label{Line:StatisticalComplexity}
		\State \Return $\text{H}\times\text{C}$
		
		\vspace{0.15cm}
		
		\Function{\texttt{WATG}}{signal.\mbox{1-D}, D, $\tau$}
		\State patterns $\gets$ \label{Line:WATGBegin} \texttt{BPSymbolization}(signal.\mbox{1-D}, D, $\tau$) \label{Line:BP}
		\State transitions $\gets$ \texttt{transitions}(patterns) \label{Line:Transitions}
		\State graph $\gets$ \texttt{weigthGraph}(signal.\mbox{1-D}, transitions) \label{Line:WG}
		\State Probs $\gets$ \texttt{as.vector}(graph) \label{Line:WATGEnd}
		\State \Return Probs
		\EndFunction
	\end{algorithmic}
\end{algorithm}

\begin{figure}
	\centering
	\includegraphics[width=.4\textwidth]{Figures/MethodologySAR.pdf}
	\caption{Outline of the methodology used for the classification of textures.}
	\label{fig:WATG}
\end{figure}

\section{Image Dataset}

We used the HH backscatter magnitudes of three quad-polarimetric L-band SAR images from the NASA Jet Propulsion Laboratory’s (JPL’s) uninhabited aerial vehicle synthetic aperture radar (UAVSAR) sensor with $L=36$ nominal looks:
\begin{itemize}
	\item forest and pasture region of Sierra del Lacandón National Park, Guatemala, (acquired on April 10, 2015)\footnote{\protect{\url{https://uavsar.jpl.nasa.gov/cgi-bin/product.pl?jobName=Lacand_30202_15043_006_150410_L090_CX_01\#dados}}}. 
	The image has $8917 \times 3300$ pixels with 
	\SI[inter-unit-product =$\cdot$]{10}{\meter} $\times$ \SI[inter-unit-product =$\cdot$]{2}{\meter} resolution.
	%
	\item ocean regions from Cape Canaveral Ocean (acquired on September 22, 2016).
	The image has $7038 \times 3300$ pixels with 
	\SI[inter-unit-product =$\cdot$]{10}{\meter} $\times$ \SI[inter-unit-product =$\cdot$]{2}{\meter} resolution;
	%
	\item urban area of the city of Munich, Germany (acquired on June 5, 2015)\footnote{\protect{\url{https://uavsar.jpl.nasa.gov/cgi-bin/product.pl?jobName=munich_19417_15088_002_150605_L090_CX_01\#data}}}.
	The image has $5773 \times 3300$ pixels with 
	\SI[inter-unit-product =$\cdot$]{10}{\meter} $\times$ \SI[inter-unit-product =$\cdot$]{3}{\meter} resolution.
\end{itemize}

We manually selected $200$ samples of size $128 \times 128$ to compose the dataset used in the experiments.
It is organized as follows:
$40$ samples from Guatemalan forest regions;
$40$ samples from Guatemalan pasture regions;
$80$ samples from the oceanic regions of Cape Canaveral, divided into two types with different contrast; and
$40$ samples of urban regions of the city of Munich.
Fig.~\ref{fig:samples&Series} shows examples of each. In our analysis, both types of ocean images are grouped.

We randomly split the samples in training (\SI{85}{\percent}) and test (\SI{15}{\percent}) sets.
We used the first set to train a $k$-nearest neighbor classifier algorithm with tenfold
cross-validation.

\begin{figure*}
	\centering
	\subfigure[Forest\label{Fig:figGuatemala}]{\includegraphics[width = 29mm]{Figures/guatemalaflorest.pdf}}
	\subfigure[Sea -- lower contrast\label{Fig:figCape1}]{\includegraphics[width = 29mm]{Figures/Cape1.pdf}}
	\subfigure[Sea -- higher contrast\label{Fig:figCape2}]{\includegraphics[width = 29mm]{Figures/Cape2.pdf}}
	\subfigure[Urban\label{Fig:figMunich}]{\includegraphics[width = 30mm]{Figures/munichUrban.pdf}}
	\subfigure[Pasture\label{Fig:figPasture}]{\includegraphics[width = 30mm]{Figures/pasture.pdf}}
	
	\subfigure[Signal -- forest\label{Fig:TSGuatemala}]{\includegraphics[width = 30mm]{Figures/tsguatemala.pdf}}
	\subfigure[Signal -- sea\label{Fig:TSCape1}]{\includegraphics[width = 30mm]{Figures/tscape1.pdf}}
	\subfigure[Signal -- sea\label{Fig:TSCape2}]{\includegraphics[width = 30mm]{Figures/tscape2.pdf}}
	\subfigure[Signal -- sea\label{Fig:TSMunich}]{\includegraphics[width = 30mm]{Figures/tsmunich.pdf}}
	\subfigure[Signal -- pasture\label{Fig:TSPasture}]{\includegraphics[width = 30mm]{Figures/tspasture.pdf}}
	\caption{Types of regions (Guatemala forest, Canaveral ocean types 1 and 2, Munich urban area, and pasture area) and their signal representation.}
	\label{fig:samples&Series}
\end{figure*}

\section{Experimental Results and Analysis}\label{Results}

In this section, we describe  the classification process, and 
the results of applying WATG.
To assess the performance of the technique here proposed, we analyze the impact of its parameters and compare its results in the classification with other methods.

\subsection{Properties}\label{Sec:WATGProperties}

We conducted two experiments to analyze the response of WATG to different noise levels and image rotations.
Our truth is the deterministic image generated by the function:
\begin{equation*}
z(x,y) = \sin (4x + 0.5y), 
\end{equation*}
where  $x, y \in [-2\pi, 2\pi]$.
Fig.~\ref{Fig:TestZ} shows this function as a $128\times128$-pixel patch.

The speckle noise was modeled as outcomes of independent, identically distributed unitary-mean Gamma random variables with shape parameter $L$ (the number of looks, which controls the signal-to-noise ratio)
$W(L) \sim \Gamma(L, L)$,
with $L \in \{1, 100, 150, 200, 250, 300, 350, 400, 450, 500\}$.
The observed images $I(L)$ are the pixelwise product of $z$ and $w(L)$.
Fig.~\ref{Fig:TestZW} shows the product $I(100)$.

\begin{figure}
	\centering
	\subfigure[$z$]{\label{Fig:TestZ}\includegraphics[width = 45mm]{Figures/z.png}}
	\subfigure[$I(100)$]{\label{Fig:TestZW}\includegraphics[width = 45mm]{Figures/z100.png}}
	\subfigure[$I^{\text{T}}(500)$]{\label{Fig:TestRot}\includegraphics[width = 45mm]{Figures/z500t.png}}
	\caption{Ground truth, speckled, and speckled transposed versions.}
	\label{Fig:TestZWL}
\end{figure}

Fig.~\ref{Fig:TestSpeckleHC} shows how the point in the $H\times C$ plane varies according to the level of noise introduced.
The ground truth (identified as ``\textsf{0}'') has relatively low entropy and is close to the maximum complexity (the continuous line is the upper bound).
This behavior is typical of deterministic sequences.
Observe that when we inject single-look speckle ($L=1$), the entropy increases along with the complexity.
Thus, the technique is able to identify the deterministic component even when it is embedded in the strongest possible speckle noise.
The point ``\textsf{1}'' shifts towards ``\textsf{0}'' when the signal-to-noise progressively increases.

\begin{figure}
	\includegraphics[width=\textwidth]{Figures/Waves1.pdf}
	\caption{Modifications to the $H \times C$ Plane features by adding different multiplicative noises.}
	\label{Fig:TestSpeckleHC}
\end{figure}

We also verified that the points in the $H\times C$ plane are almost insensitive to rotations.
Fig.~\ref{Fig:TestRot} shows the transpose of $I(500)$.
In all cases, the coordinates $(h,c)$ of the transposed noisy images were equal, up to the fourth decimal place, to those of the original version.

These experiments provide evidence that the WATG mapping is little sensitive to rotations (thanks to the use of Hilbert-Peano curves), and that it can identify the presence of underlying structural information in the presence of varying levels of speckle.


\subsection{Analysis of ordinal patterns methods}

Fig.~\ref{fig:samples&Series} shows examples of the ocean, forest, urban, and pasture both as image patches and as 1-D sequences after the linearization process.

The variation in the magnitude of the targets' backscatter and, consequently, in the intensity of the image pixels, depends on the intrinsic properties of the region under analysis.
Urban targets usually exhibit the strongest variation, followed by forest, pasture, forests, and finally, water bodies.
By adding such information related to the amplitude, the proposed method is able to increase, compared to traditional methods, the granularity of information captured by ordinal patterns.

As already described in Section~\ref{chapter:WATG}, our proposal weights the edges in terms of the difference of amplitudes.
As expected, the greatest impact is observed on the transition graphs obtained from urban areas.
The urban area \mbox{1-D} signal shown in Fig.~\ref{fig:samples&Series} has the largest dynamic range.
Fig.~\ref{fig:graphs} shows how this information alters the weights of the transition graph.
Notice, in particular, that 
$(v_{\widetilde \pi^3_{123}}, v_{\widetilde \pi^3_{123}})$ almost doubled, while 
$(v_{\widetilde \pi^3_{312}}, v_{\widetilde \pi^3_{231}})$ and $(v_{\widetilde \pi^3_{213}}, v_{\widetilde \pi^3_{132}})$ became negligible.

\begin{figure*}
	\centering
	\subfigure[Transition Graph\label{Fig:UrbanTG}]{\includegraphics[width = 75mm]{Figures/UrbanTG.pdf}}
	\subfigure[WATG\label{Fig:UrbanWATG}]{\includegraphics[width = 75mm]{Figures/UrbanWATG.pdf}}
	\caption{Difference of edges weights between the transition graph and the weighted graph of ordinal patterns transitions; urban area, with dimension 3 and delay 1.}
	\label{fig:graphs}
\end{figure*}

Fig.~\ref{fig:2DVisualization} shows the impact of using the data amplitude information on the weights of the transition graphs.
Bandt-Pompe symbolization was the first method based on ordinal patterns proposed in the literature.
As shown in Fig.~\ref{Fig:bp}, it provides limited separation of textures.
The transition graphs (Fig.~\ref{Fig:tg})improve the spread of the features, but with some amount of confusion.
Our proposal, shown in Fig.~\ref{Fig:watgD3T1}, produces well-separated features.
In this way, we were able to obtain, for this experiment, a perfect characterization and, consequently, the high descriptive power of the regions.

\begin{comment}
\begin{figure}
	\includegraphics[width=\columnwidth]{Figures/HCAnalysis.pdf}
	\caption{Location of Guatemala (forest), Cape Canaveral (ocean), and Munich (urban) in the $H \times C$ plane for dimension 3 and delay 1. 
	The continuous curves correspond to the maximum and minimum values of $C$ as a function of $H$.}
	\label{fig:plotsHC}
\end{figure}
\end{comment}

\subsection{Experiments on sliding window selection}

In this section, we analyze the parameters of the proposed method and its impact on textures classification.
McCullough et al.~\cite{McCullough2015lagged} report that inadequate values may hinder important characteristics of the phenomenon under analysis.
The two parameters of the transition graph are the dimension $D$, and the delay $\tau$.
In the experiments below, we present the results in the classification using different values of these parameters.

The classification method's performance based on ordinal patterns is sensitive to window size, the embedding dimension, and the delay.
In techniques based in Bandt-Pompe symbolization, for a fixed signal, as the size of the embedding dimension decreases, more ordinal patterns are produced.
Therefore, we acquire a higher granularity of information about the dynamics of the system and, consequently, we capture more spatial dependencies between the elements.

Fig.~\ref{fig:ROC} shows the ROC plane for different values of $D \in \{3, 4, 5, 6 \} $ and $\tau \in \{1, 2, 3, 4, 5 \}$ to select the best configuration.
The configurations that extracted most information from the \mbox{1-D} signal and, thus, that presented the best results in the experiments, are $(D = 3, \tau = 1)$ and $(D = 4, \tau = 1)$.
The technique, thus, shows its best performance choosing the parameters with the lowest computational cost.

\begin{figure}
	\includegraphics[width=\columnwidth]{Figures/ROC.pdf}
	\caption{Evaluation of the sliding window parameters using ROC curve}
	\label{fig:ROC}
\end{figure} 

Figure~\ref{fig:Regions} shows the points in the $H\times C$ produced by the same samples with all the parameters mentioned above.
The spatial distribution of the points changes with the parameters and specific configurations promote better separation.
This figure shows that the discrimination ability decreases with increasing $\tau$.
Larger values of delay dilute the spatial dependence, as neighboring points in the sample tend to be more distant in the image.
For this reason, we use $\tau=1$. Figure~\ref{fig:Regions} suggests that only one feature ($H$ or $C$) is sufficient to discriminate the classes studied. Although this is true for the experiments herein conducted, we opt to preserve the most common ordinal pattern analysis, which uses both features. As we studied only homogeneous patches, we still do not know how this approach performs with heterogeneous patches. %In this last case, it is necessary to detect borders to segment the images into homogeneous patches. 
For this last situation, we may need both features.

\begin{sidewaysfigure*}
	\centering
	\includegraphics[width=1\textwidth]{Figures/WATGHC.pdf}
	\caption{Characterization resulting in $H \times C$ Plane from the application of the Hilbert-Peano curve in WATG on textures of different regions: Guatemala (forest), Cape Canaveral (ocean) and Munich (urban). 
	The continuous curves correspond to the maximum and minimum values of $C$ as a function of $H$.}
	\label{fig:Regions}
\end{sidewaysfigure*}

Considering $\tau=1$ (first column of Fig.~\ref{fig:Regions}), 
we also notice that $D=3$ produces the best separation among classes.
Increasing $D$ also increases the Statistical Complexity; this is noticeable for the Forest class.
The other effect of considering larger values of $ D $ is an increased Entropy of Ocean and an undesirable overlap with Urban samples.

\subsection{Quantitative Evaluation}

We present a comparison between our proposal and other methods for texture characterization and classification.
We use the following ten methods: 
Gabor filters~\cite{weldon1996efficient},  
Histogram of oriented gradients (HOG)~\cite{dalal2005histograms},
Gray-level co-occurrence matrices (GLCM)~\cite{kourgli2012texture}, 
Speeded-Up Robust Features (SURF)~\cite{bay2006surf},
Short Time Fourier Transform (STFT)~\cite{portnoff1980time} with SURF,
Bandt-Pompe probability distribution~\cite{Bandt2002Permutation}, 
Ordinal patterns transition graphs~\cite{Borges2019Transition},
Weighted Permutation Entropy (WPE)~\cite{Fadlallah2013Weightedpermutation},
Fine-Grained Permutation Entropy (FGPE)~\cite{xiao2009fine} with $\alpha = 0.5$, and
Amplitude-Aware Permutation Entropy (AAPE)~\cite{azami2016amplitude} with $A = 0.5$.
As in~\cite{guan2019covariance}, 
we computed four statistics from co-occurrence matrices: contrast, correlation, energy, and homogeneity.
Likewise, we implemented the Gabor filters in five scales and eight orientations; using the energy, we obtained an $80$-dimensional feature vector for each patch.
For the HOG technique, we used image pixels divided into equal cells of $3 \times 3$ pixels, and for each cell, we computed 6-bin histograms ranging from \SIrange{0}{180}{\degree} or \SIrange{0}{360}{\degree}.

The results of visualizing the features of these descriptors can be seen in Fig.~\ref{fig:2DVisualization}.
As we can see, the algorithms that combine ordinal characteristics with amplitude information (WPE, AAPE, and FGPE) reduce the intra-class distance, but not are good discriminators.
Although the algorithms that have the largest descriptors (Gabor and HOG) do not present a good visualization when applying the dimensionality reduction technique, they presented excellent results in the quantitative evaluation.


\begin{figure*}
	\centering
	\subfigure[Gabor Filter\label{Fig:Gabor}]{\includegraphics[width = .3\columnwidth]{Figures/Gabor.pdf}}
	\subfigure[HOG\label{Fig:hog}]{\includegraphics[width = .3\columnwidth]{Figures/HOG.pdf}}
	\subfigure[GLCM\label{Fig:glcm}]{\includegraphics[width = .3\columnwidth]{Figures/GLCM.pdf}}
	\subfigure[SURF\label{Fig:surf}]{\includegraphics[width = .3\columnwidth]{Figures/Surf.pdf}}
	\subfigure[STFT + SURF\label{Fig:stft}]{\includegraphics[width = .3\columnwidth]{Figures/StftSurf.pdf}}
	\subfigure[Bandt-Pompe\label{Fig:bp}]{\includegraphics[width = .3\columnwidth]{Figures/BP.pdf}}
	\subfigure[Transition Graph\label{Fig:tg}]{\includegraphics[width = .3\columnwidth]{Figures/TG.pdf}}
	\subfigure[WPE\label{Fig:wpe}]{\includegraphics[width = .3\columnwidth]{Figures/WPE.pdf}}
	\subfigure[AAPE ($A = 0.5$)\label{Fig:aape}]{\includegraphics[width = .3\columnwidth]{Figures/AAPEA05.pdf}}
	\subfigure[FGPE ($\alpha = 0.5$)\label{Fig:fgpe}]{\includegraphics[width = .3\columnwidth]{Figures/FGPEA05.pdf}}
	\subfigure[WATG\label{Fig:watgD3T1}]{\includegraphics[width = .3\columnwidth]{Figures/WATG.pdf}}
	\caption{Bi-dimensional visualization of the descriptors obtained by our proposal and other methods of texture characterization and classification. 
	For those who provide more than two features, we apply the PCA algorithm to obtain such a visualization.}
	\label{fig:2DVisualization}
\end{figure*}

We classified the features using the $k$-nearest neighbor algorithm with Euclidean distance, selecting the value of $k$ with the automatic grid search method of the Caret R package~\cite{kuhn2008building}.
For validation, we used $10$-fold cross-validation.
More details about the classifier and the sampling can be seen in~\cite{mitchell1997machine}.

Table~\ref{tab:result1} presents the number of features each method produces, as well as its performance at classifying the $200$ samples.
We assessed the effectiveness of each approach using the following metrics. We used the first two metrics (Recall and Precision) to evaluate classifiers' per class performance and the last three metrics (Average Accuracy, Micro F1-score, and Macro F1-score) to evaluate the overall performance of the multi-class classifiers. We denote $\text{TP}_i$, $\text{TN}_i$, $\text{FP}_i$, and $\text{FN}_i$ as the true positives, true negatives, false positives, and false negatives counts of a given class $i$, among a set of $K$ classes, respectively.
\begin{itemize}
	\item Recall or True Positive Rate of the class $i$ (TPR$_i$): 
	\begin{equation*}
	\text{TPR$_i$} = \frac{\text{TP}_i}{\text{TP}_i + \text{FN}_i}.
	\end{equation*}
	%
	\item Precision or Positive Predictive Value of the class $i$ (PPV$_i$): \begin{equation*}
	\text{PPV$_i$} = \frac{\text{TP}_i}{\text{TP}_i + \text{FP}_i}.
	\end{equation*}
	%
	\item Average Accuracy (AA):
	\begin{equation*}
	\text{AA} = \sum^{\text{K}}_{i = 1}\left\{\frac{\text{TP}_i + \text{TN}_i}{\text{TP}_i + \text{TN}_i + \text{FP}_i + \text{FN}_i}\right\} 
	\end{equation*}
	%
	\item Micro F1-score: 
	\begin{equation*}
    	\text{PPV}_\mu = \frac{\sum^{\text{K}}_{i = 1}\text{TP}_i}{\sum^{\text{K}}_{i = 1}\text{TP}_i + \sum^{\text{K}}_{i = 1}\text{FP}_i}
	\end{equation*}
	\begin{equation*}
	    \text{TPR}_\mu =  \frac{\sum^{\text{K}}_{i = 1}\text{TP}_i}{\sum^{\text{K}}_{i = 1}\text{TP}_i + \sum^{\text{K}}_{i = 1}\text{FN}_i}
	\end{equation*}
	\begin{equation*}
	    \text{F1-score}_\mu = 2 \, \frac{\text{PPV}_\mu \times \text{TPR}_\mu}{\text{PPV}_\mu + \text{TPR}_\mu}.
	\end{equation*}
	%
	\item Macro F1-score: 
	\begin{equation*}
	\text{PPV}_M = \frac{1}{K} \sum^{\text{K}}_{i = 1}\frac{\text{TP}_i}{\text{TP}_i + \text{FP}_i}
	\end{equation*}
	\begin{equation*}
	\text{TPR}_M =  \frac{1}{K} \sum^{\text{K}}_{i = 1}\frac{\text{TP}_i}{\text{TP}_i + \text{FN}_i}
	\end{equation*}
	\begin{equation*}
	\text{F1-score}_M = 2 \, \frac{\text{PPV}_M \times \text{TPR}_M}{\text{PPV}_M + \text{TPR}_M}.
	\end{equation*}
\end{itemize}

\begin{table}
	\centering
	\caption{Experimental results using $k$-NN}
	\label{tab:result1}
    \scalebox{0.65}{
	\begin{tabular}{lrrrr*9{r}}
		\toprule
		\multirow{2}{*}{Method} & \multirow{2}{*}{\# features} & & TPR & & & & PPV & & & \multirow{2}{*}{AA} & \multirow{2}{*}{$\text{F1-Score}_{\mu}$} & \multirow{2}{*}{$\text{F1-Score}_M$} \\ 
		\cmidrule(lr){3-6} 
		\cmidrule(lr){7-10}
		&   & Forest & Pasture & Ocean & Urban & Forest & Pasture & Ocean & Urban & &  \\ 
		\cmidrule(lr){1-1}
		\cmidrule(lr){2-2}
		\cmidrule(lr){3-6}
		\cmidrule(lr){7-10}
		\cmidrule(lr){11-11}
		\cmidrule(lr){12-12}
		\cmidrule(lr){13-13}
		Gabor & 80 & 
		1.000 & 1.000 & 1.000 & 1.000 & 
		1.000 & 1.000 & 1.000 & 1.000 & 
		1.000 & 1.000 & 1.000\\
		HOG & 54 & 
		1.000 & 1.000 & 1.000 & 1.000 & 
		1.000 & 1.000 & 1.000 & 1.000 & 
		1.000 & 1.000 & 1.000\\
		GLCM & 32 & 
		0.833 & 1.000 & 1.000 & 0.833 &
		1.000 & 0.857 & 0.923 & 1.000 &
		0.967 & 0.980 & 0.970\\
		SURF & 1856 & 
		0.500 & 0.000 & 1.000 & 0.000 & 
		1.000 & 0.000 & 0.444 & 0.000 &
		0.467 & 0.666 & 0.572\\
		STFT + SURF  & 1856 & 
		0.166 & 0.000 & 0.833 & 0.166 & 
		0.250 & 0.000 & 0.416 & 0.500 &
		0.300 & 0.462 & 0.292\\
		\cmidrule(lr){1-1}
		\cmidrule(lr){2-2}
		\cmidrule(lr){3-6}
		\cmidrule(lr){7-10}
		\cmidrule(lr){11-11}
		\cmidrule(lr){12-12}
		\cmidrule(lr){13-13}
		Bandt-Pompe & 2 & 
		0.333 & 1.000 & 0.750 & 1.000 &
		0.500 & 0.857 & 0.750 & 0.857 &
		0.600 & 0.776 & 0.633\\ 
		Transition Graph & 2 & 
		0.833 & 0.666 & 0.833 & 1.000 &
		0.833 & 0.800 & 0.769 & 1.000 &
		0.767 & 0.929 & 0.875\\
		\cmidrule(lr){1-1}
		\cmidrule(lr){2-2}
		\cmidrule(lr){3-6}
		\cmidrule(lr){7-10}
		\cmidrule(lr){11-11}
		\cmidrule(lr){12-12}
		\cmidrule(lr){13-13}
		WPE & 2 & 
		1.000 & 0.833 & 1.000 & 0.833 &
		0.857 & 0.833 & 1.000 & 1.000 &
		0.933 & 0.868 & 0.779\\ 
		AAPE & 2 & 
		0.666 & 1.000 & 1.000 & 1.000 &
		1.000 & 0.857 & 0.923 & 1.000 &
		0.833 & 0.947 & 0.896\\ 
		FGPE & 2 & 
		0.666 & 0.666 & 1.000 & 1.000 &
		0.800 & 0.666 & 0.923 & 1.000 &
		0.767 & 0.868 & 0.711\\ 
		\cmidrule(lr){1-1}
		\cmidrule(lr){2-2}
		\cmidrule(lr){3-6}
		\cmidrule(lr){7-10}
		\cmidrule(lr){11-11}
		\cmidrule(lr){12-12}
		\cmidrule(lr){13-13}
		WATG & 2 & 
		1.000 & 1.000 & 1.000 & 1.000 &
		1.000 & 1.000 & 1.000 & 1.000 &
		1.000 & 1.000 & 1.000\\
		\bottomrule
	\end{tabular}}
\end{table}

Table~\ref{tab:result1} shows that, among the methods of weighting ordinal patterns, FGPE produced the worst results: $\text{AA} = \SI{76.7}{\percent}$, $\text{F1-score}_\mu = \SI{86.8}{\percent}$, and $\text{F1-score}_M = \SI{71.1}{\percent}$.
%
WPE also produced a low F1-score, but it produced consistently better results in the other metrics, presenting $\text{AA} = \SI{93.3}{\percent}$.
%
AAPE achieved one of the best F1-score results: $\text{AA} = \SI{83.3}{\percent}$, $\text{F1-score}_\mu = \SI{94.7}{\percent}$, and $\text{F1-score}_M = \SI{89.6}{\percent}$.
%
WATG, considering the transition graph of ordinal patterns, can better describe the textures presented, as it achieves the best performance achievable in all metrics: \SI{100}{\percent}.

\mbox{STFT + SURF} produced the worst results: $\text{AA} = \SI{30.0}{\percent}$, $\text{F1-score}_\mu = \SI{46.2}{\percent}$, and $\text{F1-score}_M = \SI{29.2}{\percent}$.
%
SURF alone provided a better performance: $\text{AA} = \SI{46.7}{\percent}$,  $\text{F1-score}_\mu = \SI{66.6}{\percent}$, and $\text{F1-score}_M = \SI{57.2}{\percent}$.
%
HOG and Gabor filters achieved the highest success rates
among all the handcrafted methods here considered: $\text{AA} = \SI{100}{\percent}$, $\text{F1-score}_\mu = \SI{100}{\percent}$, and $\text{F1-score}_M = \SI{100}{\percent}$.
However, WATG achieves that same performance using only $2$ features.
This reduction implies less computational power requirement and avoids the curse of dimensionality~\cite{TheCursesofDimensionality2018}.
Moreover, the features it is based upon are fully interpretable.

\section{Conclusions}\label{Conclusion}

We presented and assessed a new method of analysis and classification of SAR image textures.
This method consists of three steps: 
(1)~linearization, 
(2)~computing the Weighted Ordinal Pattern Transition Graph, and 
(3)~obtaining Information Theory descriptors.
A simple $k$-NN algorithm applied to the pairs Entropy-Statistical~Complexity classifies the data with \SI{100}{\percent} performance.
In addition to such perfect separation among urban, pasture, ocean, and forest areas, the proposed descriptors are interpretable in terms of the degree and structure of the spatial dependence among observations.

Experiments using patches from UAVSAR images showed that the proposal performs better than GLCM, Bandt-Pompe, Transition Graphs, SURF, \mbox{STFT + SURF}, and other techniques which also employ amplitude information in the analysis of ordinal patterns. 
Our approach provides the same quality of results obtained with Gabor filters and HOG.
However, while Gabor filters employ $80$ features and HOG uses $54$ features, our proposal requires only two.
This such reduced dimensionality consists of a huge advantage over the other techniques, with added values:
Firstly, by reducing the dimension of the features to \mbox{2-D}, we can visualize the differences between the classes of regions analyzed.
Secondly, for machine learning algorithms, the smaller the number of dimensions, the faster the training process is, and the less storage space is required.
Thirdly, we managed to avoid overfitting, a recurring problem in data with high dimensionality.

We also observed that only one feature ($H$ or $C$) is enough to discriminate the classes with the same reported performance. We opted to preserve both features because we consider that this study shed light on a novel way of SAR image analysis. Thus, we preferred to stick to the most common ordinal pattern analysis using the $H \times C $ plane. 
As future work, we consider investigating the discriminative power of these features in more complex situations.

Our approach is robust to rotations and the presence of speckle noise. The behavior showed in Figure~\ref{Fig:TestSpeckleHC} shows that our approach can capture the speckle contamination adequately.

Since the application of this work is limited to texture patches from homogeneous regions, we aim to study the possible impacts of heterogeneous areas, such as mixed culture and urban regions.
	
\section{Reproducibility and Replicability} 

Following the guidelines presented in Ref.~\cite{ABadgingSystemforReproducibilityandReplicabilityinRemoteSensingResearch}, the text, source code, and data used in this study are available at the \textit{SAR-WATG} repository \url{https://github.com/EduardaChagas/SAR-WATG}.
The information includes a link to download the $200$ labeled samples we employed in the analysis.
\chapter{HC-PCA Confidence Regions}\label{chapter:HCPCA}

Although the limits of $H \times C$ are well defined, a complete characterization of its intrinsic topology is an open problem, due to the restrictions imposed by its curvilinear space.
The lack of knowledge of the joint distribution of the points obtained by this plane, due to the existing correlation between its variables, prevents the studies on test statistics for typical time series in this characterization space.
However, with the knowledge of the expected variability of such points, according to the underlying dynamics, we can test hypotheses for a wide variety of models.
Results in this direction can be found in the literature.
\cite{RandomNumberGeneratorsCausality} showed that the Complexity-Entropy plane ($ H \times C $) is a good indicator of the results of Diehard tests for pseudo-random number generators.
\cite{De_Micco_2008} evaluated ways to improve pseudo-random sequences for their representation in this plane.

In this context, a open problem present in the characterization of sequences using the $H \times C$ plane is the absence of a representative metric distance, which makes it difficult to build confidence regions.
Thus, in the proposed approach, we opted for the construction of empirical confidence regions obtained through an orthogonal projection of data in the space of principal components.
Therefore, the larger is the data set used to build the region, the more representative it will be.

To investigate the power of representation of the proposed confidence region, we defined as a first application a review of the results obtained in the literature using the $H \times C$ plane in pseudo-random sequences.
Thus, the necessary input to our algorithm consists of sequences of true random generated by physical procedures.
As a result, we verify how the sequences previously analyzed behaved in this new analysis scenario.

Below we list the main strategies and sequences analyzed by the literature for the study of randomness using the information theory descriptors.
Soon after, the test and characterization framework is presented.

\section{Complexity-Entropy plane in the literature}

The first works on the characterization of white noises with permutation entropy arose from the need to discriminate them in relation to chaotic maps~\citep{rosso2013characterization, xiong2020complexity}.

Stationary time series can be decomposed into two main components:
\begin{enumerate}
    \item \textbf{Deterministic}: it is described by a linear combination of its own past,
    \item \textbf{Random}: it is a component of the moving average of finite order.
\end{enumerate}
In this way, chaotic systems produce sequences composed of a physical structure, easily captured by measures of complexity.
Thus, it was found that the measure of statistical complexity was able to efficiently quantify the performance of pseudorandom number generators, expanding the possibilities of using information theory descriptors with the Bandt-Pompe symbolization~\citep{larrondo2002statistical, gonzalez2005statistical, RandomNumberGeneratorsCausality}.

Table~\ref{Tab:Literature} presents a summary of the main works in the literature that perform analysis of non-chaotic algorithmic generators, according to their features in $H \times C$.
We also provide the length $T$ and embedding dimension
$D$ of the time series under scrutiny. 
The following algorithmic generators were analyzed:
\begin{itemize}
    \item Mother RNG, available in Marsaglia website~\citep{marsaglia1994yet} (MOT);
    \item Multiple with carry RNG (MWC)~\citep{marsaglia1994yet};
    \item Combo RNG (COM)~\citep{marsaglia1994yet};
    \item Lehmer RNG (LEH)~\citep{payne1969coding};
    \item Fractional Gaussian noise with $\alpha = 0$ (fGn)~\citep{bardet2003generators};
    \item $f^{-k}$ noise with $k = 0$~\citep{larrondo2012matlab};
    \item Linear Congruential Generator (LCG)~\citep{knuth1997sorting}.
\end{itemize}

\begin{table}[htb]
    \caption{Result of the main works of white noise sequences analysis in the $H \times C$ plane.}
    \label{Tab:Literature}
    \centering
    \begin{tabular}{llccccb{5.4em}}
    \toprule
	Reference & PRNG & $T$ & $D$ & $H$ & $C$ & Considered random?\\ 
	\midrule
	\citeauthor{larrondo2013statistical} (\citeyear{larrondo2013statistical}) &  MOT & NA & 6 & $\cong 0.9969$ & $\cong 0$ & no\\
	\cmidrule(lr){1-7}
	\citeauthor{gonzalez2005statistical} (\citeyear{gonzalez2005statistical})  &  MWC & 65536 & NA & $\cong 1$ & $0.3$ & yes\\
	 &  MOT & 65536 & NA & $\cong 1$ & $0.3$ & yes\\
	 &  COM & 65536 & NA & $\cong 1$ & $0.05$ & yes\\
	\cmidrule(lr){1-7}
	\citeauthor{RandomNumberGeneratorsCausality} (\citeyear{RandomNumberGeneratorsCausality}) &  LEH & \num[scientific-notation=true]{5 e6} & 5 & $\cong 1$ & $10^{-4}$ & yes\\
	 &  MOT & \num[scientific-notation=true]{5 e6} & 5 & $\cong 1$ & $10^{-4}$ & yes\\
	 &  MWC & \num[scientific-notation=true]{5 e6} & 5 & $\cong 1$ & $10^{-4}$ & yes\\
	\cmidrule(lr){1-7}
	\citeauthor{rosso2013characterization} (\citeyear{rosso2013characterization}) &  LCG & \num[scientific-notation=true]{1 e7} & 6 & $0.997871$ & $0.005101$ & no\\
	\cmidrule(lr){1-7}
	\citeauthor{xiong2020complexity} (\citeyear{xiong2020complexity}) &  fGn & \num[scientific-notation=true]{2 e17} & 6 & $\cong 1$ & $\cong 0$ & yes\\
	 & $f^{-k}$ & \num[scientific-notation=true]{2 e17} & 6 & $\cong 1$ & $\cong 0$ & yes\\
	\bottomrule
    \end{tabular}
\end{table}

None of these works provide $p$-values or hypothesis tests for their analysis. 
The focus of the articles is finding the set of descriptors that best discriminates chaos from noise. 
The authors make assessments about the randomness of sequence on ad hoc visual inspection of the point’s location in the $H \times C$ plane. 
Our work fills such a gap for finite sequences of white noise and proposes a methodology that can be extended to any other situation.

\section{Proposed Method}

In this section, we formalize the task of building a confidence region in the Entropy-Complexity manifold.
Then, we present our proposal to change space through the algorithm of the principal components analysis.
Our goal is to find a latent space representative of the data, without the restrictions of a curvilinear space.
Through this new representation of the data, we calculate empirical regions with different levels of confidence.
Finally, after calculating these regions, we build a test statistic that determines the probability that a given sequence belongs to the distribution of the points provided.

\subsection{Overall Framework}

The structure of our proposal consists of two steps:
\begin{itemize}
    \item \textbf{Empirical confidence region}: With the data present in a Euclidean plane, we can easily calculate empirical regions that involve the data with a certain level of confidence.
    \item \textbf{Construction of a test statistic}: To measure the similarity of new data sequences with the empirical points, a test statistic was proposed. 
    By acquiring a p-value less than $0.05$, we can reject the null hypothesis, which states that such data belong to the empirical probability distribution used for the construction of the confidence region.
\end{itemize}

\subsection{Empirical Confidence Regions and \texorpdfstring{$p$}--values}\label{confidenceRegions}

Our first approaches to analyzing sequences of points in the $H\times C$ plane produced by TRWNS verified that they, and usual transformations, are far from bivariate Gaussian and generalized Hyperbolic distributions \citep{MultivariateDistributionModelswithGeneralizedHyperbolicMargins}.
Different types of regression models of $C$ explained by $H$ did not produce acceptable results.
Thus, we adopted a non-parametric approach and made an empirical analysis of the data obtained from physical sources for using them as our reference in the search for confidence regions and $p$-values.

Let $\utilde{x} =(x_1, x_2, \dots, x_N)$ be $N$ times series of length $T$, and define an embedding dimension $D$.
In the sequel, whenever possible, we will omit $T$ and $D$.
For each $n=1,2,\dots, N$, the time series $x_n$ is mapped onto the point $(h_n,c_n)$ in the $H\times C$ plane, thus $\utilde{hc}=\big((h_1,c_n), (h_2,c_2), \dots, (h_N,c_N)\big)$ are the points that correspond to the $N$ time series.
Fig.~\ref{fig:Methodology1} illustrates this step.
We will obtain confidence regions and $p$-values from $\utilde{hc}$.

The first step consists in finding and applying the principal components transformation to $\utilde{hc}$.
With this, we obtain the set of uncorrelated points $\utilde{uv}=\big((u_1,v_1), (u_2,v_2),\dots,(u_N,v_N)\big)$, in which $u_n$ and $v_n$ are the first and second principal components of $h_n$ and $v_n$, respectively.
This projection allows us to obtain a ``central'' point of the data set, around which we will build a rectangular box containing \SI{100}{\minusalphapercent} of the observations, where $\alpha$ is the significance level analyzed.
Such box is a variation of the bagplot \citep{TheBagplotaBivariateBoxplot}.
Notice that finding the smallest box that encloses $k$ out of $N$ points is difficult; cf.\ the work by \citet{SmallestKEnclosingRectangleRevisited}.

For simplicity, and without loss of generality, assume $N$ is odd.
\begin{enumerate}
	\item Find the indexes that sort the values of the first principal component $\bm u=(u_1,u_2,\dots,u_N)$ in ascending order: $\bm r=(r_1,r_2,\dots,r_N)$, i.e., $u_{r_1}$ is the minimum value, and $u_{r_N}$ is the maximum value.
	\item\label{item:Median} Find the point $(u,v)$ whose first principal component is the median: $(u_{r_{(N+1)/2}}, \cdot)$. Apply the inverse principal components transformation, and obtain $\bm P'=(h',v')$. Call the corresponding time series ``emblematic time series.''
	\item\label{item:Point1} Find the point $(u,v)$ whose first principal component is the quantile $\alpha/2$: $(u_{r_{[N\alpha/2]}}, \cdot)$.
	\item\label{item:Point2} Find the point $(u,v)$ whose first principal component is the quantile $1-\alpha/2$: $(u_{r_{[N(1-\alpha/2)]}}, \cdot)$.
	\item\label{item:Point3} The values $u_{r_{[N\alpha/2]}}$ and $u_{r_{[N(1-\alpha/2)]}}$ are the rightmost and leftmost bounds of the box, respectively.
	\item\label{item:Point4} The bottom bound of the box is the smallest second principal component value whose first principal component is at least $u_{r_{[N\alpha/2]}}$; denote this values $v_{\min}$.
	\item\label{item:Point5} The top bound of the box is the largest second principal value whose first principal component is at most $u_{r_{[N(1-\alpha/2)]}}$; denote this value $v_{\max}$.
	\item The corners of the box are 
	$(u_{r_{[N\alpha/2]}}, v_{\min})$, 
	$(u_{r_{[N\alpha/2]}}, v_{\max})$, 
	$(u_{r_{[N(1-\alpha/2)]}}, v_{\min})$ and 
	$(u_{r_{[N(1-\alpha/2)]}},v_{\max})$.
	\item\label{item:BoxHxC} Apply the inverse principal components transformation to these corners obtaining $\bm P_1=(h_{v_1}, c_{v_1})$, $\bm P_2=(h_{v_2},h_{v_2})$, $\bm P_3=(h_{v_3}, c_{v_3})$ and $\bm P_4=(h_{v_4},c_{v_4})$.
\end{enumerate}

Fig.~\ref{fig:methodology} illustrates these steps.
%
Fig.~\ref{fig:Methodology1} shows the points produced by TRWNS in the $H\times C$ plane.
The blue box includes a certain percentage of points, with sides parallel to the $H$ and $C$ axes.
The area in the $H\times C$ plane overestimates the desired proportion and may include ``unacceptable'' points.
%
Fig.~\ref{fig:Methodology2} shows the previous points projected onto the principal components space (steps~\ref{item:Median} to~\ref{item:Point5}).
The red box includes the same percentage of desired points, with axes parallel to the first and second principal components.
We highlighted in red the point whose first principal component is the median of the observed values.
%
Fig.~\ref{fig:Methodology3} shows the result of projecting back the red box from the principal components space to the $H\times C$ plane (step~\ref{item:BoxHxC}).
The comparison of the red and blue boxes shows that the area has been reduced, thus improving the test's power.

\begin{figure*}
	\centering
	\subfigure[Mapping true white noise random sequences onto the $H\times C$ plane.\label{fig:Methodology1}]{\includegraphics[width=.32\linewidth]{Figures/step1}}
	\subfigure[Transformation of the points in the $H\times C$ plane by Principal Components, and determination of minimal boxes.\label{fig:Methodology2}]{\includegraphics[width=.32\linewidth]{Figures/step2}}
	\subfigure[Inverse transformation from the Principal Components plane to the $H\times C$ plane.\label{fig:Methodology3}]{\includegraphics[width=.32\linewidth]{Figures/step3}}
	\caption{Outline of the methodology used for the construction of the confidence regions.}
	\label{fig:methodology}
\end{figure*}

Algorithm~\ref{Algo:ConfidenceRegions} provides details on how we obtain the confidence regions, defined by a set of points $\bm P_1, \bm P_2, \bm P_3, \bm P_4$, for each $D\in \mathcal D$, each $T\in \mathcal T$, and each significance level $\alpha$.
We also obtain the ``emblematic point'' $\bm P'$, a kind of median point in the $H\times C$ plane for each situation.

  \begin{center}
    \scalebox{0.85}{
    \begin{minipage}{\linewidth}
\begin{algorithm}[H]
	\SetKwInOut{Input}{input}\SetKwInOut{Output}{output}
	\Input{A data base of true random values}
	\Input{The desired values of embedding dimension $\mathcal D$, sequence length $\mathcal T$, and confidence levels $\mathcal A$}
	\Output{Confidence regions as points in the $H\times C$ plane}
	\For{each $D\in \mathcal D$}{
		\For{each $T\in \mathcal T$}{
			\For{each $n=1,2,\dots,N$}{
				build the time series $p_n$ with unused values from the data base\;
				compute the point $(h_n,c_n)$ in the $H\times C$ plane that corresponds to $p_n$\; 
			}
			obtain $\text{PC}(D,T)$, the principal components transformation based on the points $(h_1,c_1),(h_2,c_2),\dots,(h_N,c_N)$, and its inverse $\text{PC}^{-1}(D,T)$\;
			%
			apply $\text{PC}(D,T)$ to the points $(h_1,c_1),(h_2,c_2),\dots,(h_N,c_N)$, and obtain $(u_1,v_1),(u_2,v_2),\dots,(u_N,v_N)$\;
			%
			find the indexes $\bm r=(r_1,r_2,\dots,r_N)$ that sort the values of the first principal component $\bm u=(u_1,u_2,\dots,u_N)$ in ascending order\;
			%
			find the point $(u,v)$ whose first principal component is the median: $(u_{r_{(N+1)/2}}, \cdot)$\;
			%
			apply the inverse principal components transformation $\text{PC}^{-1}(D,T)$ to $(u,v)$, and obtain $\bm P'=(h',v')$; call the corresponding time series ``emblematic time series''\;
			\Return{$\bm P'$}\;
			%
			\For{each confidence level $\alpha \in \mathcal A$}{
				find the point $(u,v)$ whose first principal component is the quantile $\alpha/2$: $(u_{r_{[N\alpha/2]}}, \cdot)$\;
				%
				find the point $(u,v)$ whose first principal component is the quantile $1-\alpha/2$: $(u_{r_{[N(1-\alpha/2)]}}, \cdot)$\;
				%
				the values $u_{r_{[N\alpha/2]}}$ and $u_{r_{[N(1-\alpha/2)]}}$ are the rightmost and leftmost bounds of the box, respectively\;
				%
				the bottom bound of the box is the smallest second principal component value whose first principal component is at least $u_{r_{[N\alpha/2]}}$; denote this value $v_{\min}$\;
				%
				the top bound of the box is the largest second principal value whose first principal component is at most $u_{r_{[N(1-\alpha/2)]}}$; denote this value $v_{\max}$\;
				%
				the corners of the box are 
				$(u_{r_{[N\alpha/2]}}, v_{\min})$, 
				$(u_{r_{[N\alpha/2]}}, v_{\max})$, 
				$(u_{r_{[N(1-\alpha/2)]}}, v_{\min})$ and 
				$(u_{r_{[N(1-\alpha/2)]}},v_{\max})$\;
				%
				apply the inverse principal components transformation $\text{PC}^{-1}(D,T)$ to these corners obtaining $\bm P_1=(h_{v_1}, c_{v_1})$, $\bm P_2=(h_{v_2},h_{v_2})$, $\bm P_3=(h_{v_3}, c_{v_3})$ and $\bm P_4=(h_{v_4},c_{v_4})$\;
				\Return{$\bm P_1$, $\bm P_2$, $\bm P_3$, $\bm P_4$}\;
			}
		}
	}
	\caption{Determination of confidence regions and emblematic time series}
	\label{Algo:ConfidenceRegions}
      \end{algorithm}
    \end{minipage}%
    }
  \end{center}

These confidence regions obtained provide a powerful tool to make binary assessments about the adequacy of a given time series $\bm x$ to the null hypothesis $\mathcal H_0$ that it is white noise.
More generally, we are interested in obtaining the $p$-value of $\bm x$ under $\mathcal H_0$.
We present a procedure to obtain an approximate $p$-value based on the evidence collected to build the confidence regions.

The procedure operates on the principal components space and consists of measuring the closeness between the ``emblematic point'' and the observed point.
We are given a time series $\bm x$ of size $T$, and we want its $p$-value when contrasted with TWNRS of the same size at embedding dimension $D$.
We use $N$ TWNRS of size $T$, compute their points in the $H\times C$ plane, and project them to the corresponding principal components space.
We then do the same with $\bm x$, and obtain a new point $(u_x,v_x)$.
The closer $\bm x$ is to the emblematic time series, the larger its $p$-value.
Assume that the emblematic time series is represented by $(u,v)$ in the principal components space.
We measure this closeness by building a box around $(u_x,v_x)$ that contains $(u,v)$; assume that $u_x>u$, then:
\begin{enumerate}
	\item the right side of the box is the smallest $u_j$ which is larger that $u_x$; assume it corresponds to the quantile $\eta_u$ of $\utilde u = (u_1,u_2,\dots, u_N)$. By definition, $\eta_u\geq 1/2$.
	%
	\item the left side of the box is the $1-\eta_u$ quantile of $\utilde u$.
	%
	\item the top side of the box is the smallest $v_j$ which is larger that $v_x$; assume it corresponds to the quantile $\eta_v$ of $\utilde v = (v_1,v_2,\dots, v_N)$. By definition, $\eta_v\geq 1/2$.
	%
	\item the bottom side of the box is the $1-\eta_v$ quantile of $\utilde v$.
\end{enumerate}
Fig.~\ref{fig:methodologyPvalue} illustrates theses steps.

The definition of the box for the case $u_x<u$ follows naturally, and is described in Algorithm~\ref{Algo:p-value}.
With this approach, we obtain the smallest box that (i)~contains the new point, and (ii)~is defined by observed points from TRWNS.

Such boxes are less prone to distortions in this space since the distribution of the points becomes less asymmetric than in the $H\times C$ plane; cf.\ Fig.~\ref{fig:HC-PCA}.
Algorithm~\ref{Algo:p-value} shows the details.

\begin{figure*}
	\centering
	\subfigure[Points produced by TRWNSs in the space of principal components.\label{fig:v1}]{\includegraphics[width=.325\linewidth]{Figures/PvalueStep1}}
	\subfigure[Points whose first principal component is the median (red) and quantiles of order $\eta_u$ and $1-\eta_u$ (green).\label{fig:methodologyPvalue2}]{\includegraphics[width=.325\linewidth]{Figures/PvalueStep2}}
	\subfigure[The $p$-value of the new point is the proportion of points outside the box.\label{fig:methodologyPvalue3}]{\includegraphics[width=.325\linewidth]{Figures/PvalueStep3}}
	\caption{Outline of the methodology used to calculate the $p$-value. The new point is denoted as a crossed circle.}
	\label{fig:methodologyPvalue}
\end{figure*}

\begin{algorithm}[hbt]
	\SetKwInOut{Input}{input}\SetKwInOut{Output}{output}
	\Input{The sequence $\bm x$ of length $T$ to be contrasted to the null hypothesis $\mathcal H_0$ that it is adherent to white noise}
	\Input{The embedding dimension $D$}
	\Input{$N$ points $(h_1,c_1),(h_2,c_2),\dots,(h_N,c_N)$ of true white noise series of length $T$;
		the principal components tranformation $\text{PC}(D,T)$ induced by these points;	
		the points projected onto the $H\times C$ plane: $(u_1,v_1),(u_2,v_2),\dots,(u_N,v_N)$}
	\Output{An approximate $p$-value}
	find the point $(u,v)$ whose first principal component is the median: $(u_{r_{(N+1)/2}}, \cdot)$\;
	find the point $(h,c)$ of the sequence $\bm x$\;
	find the projection $(u_x,v_x)$ of $(h,c)$ onto the principal components space using $\text{PC}(D,T)$\;	
	\eIf{$u_x > u$}
	{ $u_{r_{[N(1-\alpha/2)]}}$ is defined as the smallest element larger than $u_x$\;
		$u_{r_{[N\alpha/2]}} \leftarrow 2u - u_{r_{[N(1-\alpha/2)]}}$\;
	}
	{
		\eIf{$u_x < u$}
		{ $u_{r_{[N\alpha/2]}}$ is the largest minor element of $u_x$\;
			$u_{r_{[N(1-\alpha/2)]}} \leftarrow 2u - u_{r_{[N\alpha/2]}}$\;
		}
		{ $u_{r_{[N\alpha/2]}}$ and $u_{r_{[N(1-\alpha/2)]}}$ is equal to $u$, the median point of the first principal component\; 
		}
	}
	obtain the maximum values of the second component whose values of the first principal component are at least $u_{r_{[N\alpha/2]}}$ and at most $u_{r_{[N(1-\alpha/2)]}}$ and denote it $v_{\max}$\;
	obtain the minimum values of the second component whose values of the first principal component are at least $u_{r_{[N\alpha/2]}}$ and at most $u_{r_{[N(1-\alpha/2)]}}$ and denote it $v_{\min}$\;
	the corners of the box $b_{\alpha}(h,c)$ are 
	$(u_{r_{[N\alpha/2]}}, v_{\min})$, 
	$(u_{r_{[N\alpha/2]}}, v_{\max})$, 
	$(u_{r_{[N(1-\alpha/2)]}}, v_{\min})$ and 
	$(u_{r_{[N(1-\alpha/2)]}}, v_{\max})$\;
	%find $b_{T,D}(h,c)$: the smallest box centered at $(h',c')$ which contains $\text{PC}(D,T)(h,c)$\;
	count $n_{\bm x}$, the number of points out of the $N$ points which belong to $b_{\alpha}(h,c)$\;
	\Return{$1-n_{\bm x}/N$}
	\caption{Determination of the $p$-value of the sequence $\bm x$ under $\mathcal H_0$}\label{Algo:p-value}
\end{algorithm}


\chapter{Case study: White Noise Confidence Regions}\label{chapter:confidenceRegions}

In this section, we will present the results obtained by the first proposal for the construction of empirical confidence regions in the Complexity-Entropy plane for white noise models.
Through these regions, we want to check if the randomness present in such PRNGs can be captured by the descriptors even with the use of short sequences.
As a consequence, we observed that the proposed methodology proved to be consistent and coherent, managing to capture the randomness of truly random sequences, reproducing the results obtained by generators previously analyzed in the literature, and proving to be robust in the addition of correlation structures.

\section{Introduction}\label{conf:intro}

Several works have used deterministic and pseudorandom sequences aiming at understanding the properties of the points they produce in the $H\times C$ plane.
\cite{GeneralizedStatisticalComplexityMeasuresGeometricalAnalyticalProperties} analyzed the logistic chaotic map and discuss the boundaries of the $H \times C$ plane.
\cite{De_Micco_2009} studied chaotic components in pseudorandom number generators.
\cite{DistinguishingNoiseFromChaos}  tackled the often hard problem of distinguishing chaos from noise.
\cite{DistinguishingChaoticStochasticDynamicsTimeSeriesMultiscaleSymbolicApproach} used a multi-scale approach to analyze the interplay between chaotic and stochastic dynamics.

%With the knowledge of the expected variability of such points, according to the underlying dynamics, we can make hypothesis tests for a wide variety of models.
%Results in this direction can be found in the literature.
%\cite{RandomNumberGeneratorsCausality} showed that the Entropy-Complexity plane ($H\times C$) is a good indicator of the results of Diehard tests for pseudorandom number generators.
%\cite{De_Micco_2008} assessed ways of improving pseudorandom sequences by their representation in this plane.

However, we were able to verify that the reported works carry out statistical studies using as a basis the descriptors of Information Theory.
Motivated by previous works, in this work, we advance the state-of-the-art providing the first test for white
noise points in the $H \times C$ plane. 
In this proposal, the input is a sequence of true random observations generated by a physical-based procedure. 
We obtain the confidence regions by performing an orthogonal projection of the data onto the space of principal components, thus eliminating the restrictions imposed by the bounded space of the Complexity-Entropy plane.
Our contributions can be summarized as follows:
\begin{itemize}
    \item We provide the first contribution in constructing a test in the Complexity-Entropy Plane: we provide confidence regions and $p$-values.
    \item We evaluate this test’s size by analyzing random sequences generated by physical procedures and pseudorandom generators (PRNGs).
    \item We verify the test’s power contrasting correlated noise time series.
\end{itemize}

\section{Experimental Settings}

We evaluated the performance of the proposed method in relation to a large set of random sequences provided by state-of-the-art pseudo-random number generators.
In this section, we present the settings of the parameters that we use as a reference, the true random physical generators used to calculate the empirical distribution, and descriptive analysis of representative points in relation to the confidence regions.

\subsection{Parameters Settings and Dataset}

We conducted an ablation study to identify the influence of the parameters $T$, $D$, and $\tau$ in the construction of empirical confidence regions.
We verified that the results involving the time delay parameter variation did not show significant differences in repeated experiments; therefore, in the sequel, we did no consider $\tau$ as a determining factor.
On the other hand, we found two relevant variables: 
the length of the sequence 
and the embedding dimension.
We, thus, employed the following factors:
\begin{itemize}
\item Sequence length $T\in\mathcal T=\{T = \num[scientific-notation=true]{e3}, \num[scientific-notation=true]{5 e4}\}$,
\item Embedding dimension $D\in\mathcal D=\{3, 4, 5, 6\}$.
\end{itemize}
and kept $\tau=1$, which is the most frequently used option.
The values of $D$ are within the range recommended in the literature~\citep{PermutationEntropyBandtPompe}.

Using this parametric space, we analyzed the different degrees of information captured by the ordinal patterns formed.
For the construction of the confidence regions presented, we used:
\begin{itemize}
	\item A set of \num{104596} points in the $H \times C$ plane, corresponding to sequences of length $T = 1000$, for each value of $D\in \mathcal D$, and
	\item  a set of \num{2093} points in the $H \times C$ plane, corresponding to sequences of length $T = 50000$, for each value of $D \in \mathcal D$.
\end{itemize}

We used the \texttt R platform \citep[][v.~4.0.3]{Rmanual} for data generation and analyses, and the \texttt{ggplot2} library \cite{ggplot2Wickman} for generating the plots.

\subsection{True Random Numbers}

Random numbers are used in many fields, from gambling to criptography, aiming to guarantee a secure, realistic or unpredictable behavior. 
Pseudo randomic results can be achieved by software in a deterministic way. 
But, some applications need actual random numbers (despite the somewhat elusive nature of actual randomness).
Randomness can be observed in unpredictable real world phenomena like cathodic radiation or atmospheric noise.
%In this study we used two sources of real random numbers. 
%The first is based on vacuum states to generate random quantum numbers described by \citeauthor{RNGVacuumStates}~(\citeyear{RNGVacuumStates}), the second one is based on atmospheric noise captured by a cheap radio receiver presented at \url{www.random.org}.

In this study we used two sources of random numbers, here called true random, both from physical phenomena observation and measurement.
The first is based on vacuum states to generate random quantum numbers, the setup consists of an ordinary laser source to generate a local oscillator (LO), a half-wave plate, a polarizing beamsplitter (BPS), and two balanced detectors working together adding or subtracting the photocurrents results in a quadrature measurement of the LO or vacuum state. 
The probability distribution of the vacuum state is binned into $2^n$ equal parts (bins of same size), than, assigning a fixed bit combination of length \textit{n} to each sample point in a given bin \citeauthor{RNGVacuumStates}~(\citeyear{RNGVacuumStates}). 
The second one is based on atmospheric noise captured by a cheap radio receiver, started as a gambling engine, the randomness comes from an ordinary radio receiver that has no filter for static unwanted sounds caused by atmospheric noise, but perfect for random purposes, developed over a distributed setup with some radios located at different geographical locations sending random bits to a cloud server who process data and hosts random.org, the history, and some other information could be found at \cite{RandomOrg}.
We used \SI{54e6}{4\byte} words from each physical generator, which approximately amounts \SI{200}{\mega\byte} of data.

In Fig.~\ref{fig:white-noise}, we show how the samples used to build the empirical confidence regions are arranged in the $H \times C$ plane.


\begin{figure}
    \centering
    \includegraphics[width=\linewidth]{Figures/Points-PDF.png}
    \caption{White noise samples considered during the construction of the proposed confidence regions.}
    \label{fig:white-noise}
\end{figure}

\section{Experimental Results and Analysis}\label{Sec:Results}

\subsection{Descriptive analysis of empirical confidence regions}

The regions used as a reference in this work are obtained through true random sequences, where we extract the empirical distribution of white noises in the Complexity-Entropy plane.
Tables~\ref{Tab:Points1k} and~\ref{Tab:Points50k} list the coordinates in the $H\times C$ plane
of the emblematic point $\bm P'$ and of the four points $\bm P_1,\bm P_2,\bm P_3,\bm P_4$ that define the confidence regions at \SI{90}{\percent}, \SI{95}{\percent}, \SI{99}{\percent}, and \SI{99.9}{\percent}, 
for $D=3,4,5,6$ and $N=1000,50000$, .
The points are presented counterclockwise, starting with the one with the largest complexity.

Fig.~\ref{fig:HC-PCA} shows the results for $T = 50000$ and $D = 3,6$ in the new principal components space, along with the quantiles of order $\SI{90}{\percent}$, $\SI{95}{\percent}$, $\SI{99}{\percent}$, and $\SI{99.9}{\percent}$.
We also show the projection of the $H \times C$ plane boundaries in this space, as well the median of each data set, the latter being represented as red dots.
The confidence regions exceed the $H \times C$ boundaries, but this issue does not compromise the test's size since no points can be observed outside such boundaries.

Fig.~\ref{fig:HC-PCA} also shows that the data are not evenly distributed among the axes of the first principal component.
They tend to concentrate close to the point that corresponds to $(1,0)$ in the $H\times C$ plane.
As we use order statistics to define the confidence regions, this issue is also of little relevance for our results.
Moreover, Fig.~\ref{fig:PCA-Hist} shows that such asymmetry diminishes when the embedding dimension $D$ increases.

\begin{figure*}
	\centering
	\subfigure[Points in the Principal Components plane for for $D=3$]{\includegraphics[width=.49\linewidth]{Figures/HC-PCA-Trozos-D3N50k}}
	\subfigure[Points in the Principal Components plane for for $D=6$]{\includegraphics[width=.49\linewidth]{Figures/HC-PCA-Trozos-D6N50k}}
	\caption{Representation of true random white noise sequences of length $T = 50000$ in the PCA space for $D = 3$ and $D = 6$, and the quantiles of $\SI{90}{\percent}$, $\SI{95}{\percent}$, $\SI{99}{\percent}$, and $\SI{99.9}{\percent}$.}
	\label{fig:HC-PCA}
\end{figure*} 

As we can see in Fig.~\ref{fig:PCA-Hist} in the new representation space produced by the PCA, the data are not evenly distributed among the axes of the first main component, maintaining the character presented in the $H \times C$ plane, since such points tend to be concentrated close to the point $(1, 0)$.

\begin{figure}
	\centering
	\includegraphics[width=\linewidth]{Figures/PCA-hist-50k}
	\caption{Histograms of the first principal component for $D=3,4,5,6$}
	\label{fig:PCA-Hist}
\end{figure}

\begin{table}[hbt]
	\caption{Coordinates in the $H\times C$ plane of the emblematic series and the points that define the confidence regions at \SI{90}{\percent}, \SI{95}{\percent}, \SI{99}{\percent}, and \SI{99.9}{\percent} for $D=3,4,5,6$ and $N=1000$}
	\label{Tab:Points1k}
	\centering
    \scalebox{0.78}{
	\begin{tabular}{rccccc}
		\toprule
		& & \multicolumn{4}{c}{$N = 1000$}\\ 
		\cmidrule(lr){3-6} 
		$D$ & Point & \SI{90}{\percent} & \SI{95}{\percent} & \SI{99}{\percent} & \SI{99.9}{\percent}\\ 
		\cmidrule(lr){3-3} 
		\cmidrule(lr){4-4} 
		\cmidrule(lr){5-5} 
		\cmidrule(lr){6-6} 
		$3$ & $\bm P'$ & \multicolumn{4}{c}{$(0.9992089,0.0007800)$}\\
		& $\bm P_1$ & $(0.9973334, 0.0025601)$ & $(0.9967311, 0.0031343)$ & $(0.9953009, 0.0045054)$ & $(0.9931825, 0.0065387)$\\
		& $\bm P_2$ & $(0.9974047, 0.0026304)$ & $(0.9968219, 0.0032238)$ & $(0.9954349, 0.0046375)$ & $(0.9933704, 0.006724)$\\
		& $\bm P_3$ & $(0.9999497, 0)$ & $(0.9999398, 0)$ & $(0.9999203, 0)$ & $(0.9998925, 0)$\\
		& $\bm P_4$ & $(1, \num[scientific-notation=true]{5.17e-05})$ & $(1, \num[scientific-notation=true]{6.12e-05})$ & $(1, \num[scientific-notation=true]{8.45e-05})$ & $(1, 0.0001104)$\\ 
		\midrule
		$4$ & $\bm P'$ & \multicolumn{4}{c}{$(0.9967032, 0.0043297)$} \\
		& $\bm P_1$ & $(0.994364, 0.0081246)$ & $(0.9937138, 0.0089796)$ & $(0.,9922575 0.0108947)$ & $(0.9902578, 0.0135243)$\\
		& $\bm P_2$ & $(0.9939234, 0.0075452)$ & $(0.9932534, 0.0083741)$ & $(0.9917308, 0.0102022)$ & $(0.9897312, 0.0128318)$\\
		& $\bm P_3$ & $(0.9994791, 0.0013982)$ & $(0.9991609, 0.0018166)$ & $(0.9987924, 0.0023012)$ & $(0.9985727, 0.0025901)$\\
		& $\bm P_4$ & $(0.9990385, 0.0008188)$ & $(0.9987005, 0.0012111)$ & $(0.9982658, 0.0016087)$ & $(0.9980461, 0.0018976)$\\ 
		\midrule
		$5$ & $\bm P'$ & \multicolumn{4}{c}{$(0.9864873, 0.0245632)$}\\
		& $\bm P_1$ & $(0.9811818, 0.0321294)$ & $(0.9801289, 0.0340045)$ & $(0.977917, 0.0377295)$ & $(0.9753326, 0.0425299)$\\
		& $\bm P_2$ & $(0.9827429, 0.0350291)$ & $(0.9817117, 0.0369446)$ & $(0.9796031, 0.0408613)$ & $(0.9770187, 0.0456617)$\\
		& $\bm P_3$ & $(0.9919707, 0.0120896)$ & $(0.9909376, 0.0139279)$ & $(0.9898161, 0.0156277)$ & $(0.9892599, 0.0166608)$\\
		& $\bm P_4$ & $(0.9935319, 0.0149893)$ & $(0.9925204, 0.016868)$ & $(0.9915021, 0.0187595)$ & $(0.9909459, 0.0197926)$\\ 
		\midrule
		$6$ & $\bm P'$ & \multicolumn{4}{c}{$(0,9296429, 0.1841438)$}\\
		& $\bm P_1$ & $(0.9121895, 0.2201993)$ & $(0.9105951, 0.2239294)$ & $(0.9105951, 0.2239294)$ & $(0.9077672, 0.2305874)$\\
		& $\bm P_2$ & $(0.9146048, 0.2260776)$ & $(0.9130413, 0.2298829)$ & $(0.9130413, 0.2298829)$ & $(0.9102595, 0.2366531)$\\
		& $\bm P_3$ & $(0.9443868, 0.1418373)$ & $(0.9419202, 0.1476904)$ & $(0.9396577, 0.1531967)$ & $(0.9383611, 0.1561279)$\\
		& $\bm P_4$ & $(0.9468021, 0.1477156)$ & $(0.9443663, 0.1536439)$ & $(0.9421039, 0.1591502)$ & $(0.9408534, 0.1621937)$\\ 
		\bottomrule
	\end{tabular}}
\end{table}

\begin{table}[hbt]
	\caption{Coordinates in the $H\times C$ plane of the emblematic series and the points that define the confidence regions at \SI{90}{\percent}, \SI{95}{\percent}, \SI{99}{\percent}, and \SI{99.9}{\percent} for $D=3,4,5,6$ and $N=50000$}\label{Tab:Points50k}
	\centering
    \scalebox{0.73}{
	\begin{tabular}{rccccc}
		\toprule
		%& & \multicolumn{4}{c}{$N$}\\ 
		%\cmidrule(lr){3-6}
		& & \multicolumn{4}{c}{$N = 50000$}\\ 
		\cmidrule(lr){3-6} 
		$D$ & Point & \SI{90}{\percent} & \SI{95}{\percent} & \SI{99}{\percent} & \SI{99.9}{\percent}\\ 
		\cmidrule(lr){3-3} 
		\cmidrule(lr){4-4} 
		\cmidrule(lr){5-5} 
		\cmidrule(lr){6-6} 
		$3$ & $\bm P'$ & \multicolumn{4}{c}{$(0.9999853, \num[scientific-notation=true]{1.45e-05})$}\\
		& $\bm P_1$ & $(0.9999489, \num[scientific-notation=true]{5.06e-05})$ & $(0.9999384, \num[scientific-notation=true]{6.11e-05})$ & $(0.9999079, \num[scientific-notation=true]{9.11e-05})$ & $(0.9998625, 0.0001361)$\\
		& $\bm P_2$ & $(0.9999487, \num[scientific-notation=true]{5.04e-05})$ & $(0.9999382, \num[scientific-notation=true]{6.09e-05})$ & $(0.9999077, \num[scientific-notation=true]{9.09e-05})$ & $(0.9998622, 0.0001358)$\\
		& $\bm P_3$ & $(0.9999998, \num[scientific-notation=true]{4e-07})$ & $(0.9999994, \num[scientific-notation=true]{9e-07})$ & $(0.9999982, \num[scientific-notation=true]{2e-06})$ & $(0.9999973, \num[scientific-notation=true]{3e-06})$\\
		& $\bm P_4$ & $(0.9999996, \num[scientific-notation=true]{2e-07})$ & $(0.9999991, \num[scientific-notation=true]{7e-07})$ & $(0.999998, \num[scientific-notation=true]{1.8e-06})$ & $(0.999997, \num[scientific-notation=true]{2.7e-06})$\\ \midrule
		$4$ & $\bm P'$ & \multicolumn{4}{c}{$(0.9999394, \num[scientific-notation=true]{7.94e-05})$}\\
		& $\bm P_1$ & $(0.9999684, \num[scientific-notation=true]{3.98e-05})$ & $(0.9999725, \num[scientific-notation=true]{3.44e-05})$ & $(0.9999783, \num[scientific-notation=true]{2.68e-05})$ & $(0.9999833, \num[scientific-notation=true]{2.02e-05})$ \\
		& $\bm P_2$ & $(0.9999696, \num[scientific-notation=true]{4.13e-05})$ & $(0.9999737, \num[scientific-notation=true]{3.6e-05})$ & $(0.9999795, \num[scientific-notation=true]{2.83e-05})$ & $(0.9999845, \num[scientific-notation=true]{2.18e-05})$ \\
		& $\bm P_3$ & $(0.9998075, 0.0002508)$ & $(0.9998506, 0.0001942)$ & $(0.9998756, 0.0001615)$ & $(0.9998889, 0.000144)$ \\
		& $\bm P_4$ & $(0.9998087, 0.0002524)$ & $(0.9998518, 0.0001958)$ & $(0.9998768, 0.000163)$ & $(0.9998901, 0.0001456)$ \\ \midrule
		$5$ & $\bm P'$ & \multicolumn{4}{c}{$(0.9997616, 0.0004264)$}\\
		& $\bm P_1$ & $(0.9998172, 0.0003232)$ & $(0.9998259, 0.0003075)$ & $(0.9998428, 0.0002774)$ & $(0.9998573, 0.0002517)$ \\
		& $\bm P_2$ & $(0.9998194, 0.0003273)$ & $(0.9998282, 0.0003116)$ & $(0.999845, 0.0002814)$ & $(0.9998593, 0.0002553)$ \\
		& $\bm P_3$ & $(0.9994812, 0.0009246)$ & $(0.9996371, 0.0006455)$ & $(0.9996703, 0.0005862)$ & $(0.9996884, 0.000554)$ \\
		& $\bm P_4$ & $(0.9994834, 0.0009286)$ & $(0.9996394, 0.0006495)$ & $(0.9996725, 0.0005901)$ & $(0.9996904, 0.0005576)$ \\ \midrule
		$6$ & $\bm P'$ & \multicolumn{4}{c}{$(0.9989108, 0.0026093)$}\\
		& $\bm P_1$ & $(0.9990169, 0.002336)$ & $(0.9990368, 0.002288)$ & $(0.9990736, 0.0021997)$ & $(0.9991069, 0.0021197)$ \\
		& $\bm P_2$ & $(0.9990249, 0.0023554)$ & $(0.9990449, 0.0023074)$ & $(0.9990817, 0.0022191)$ & $(0.999115, 0.0021392)$ \\
		& $\bm P_3$ & $(0.9978983, 0.0050219)$ & $(0.998714, 0.0030633)$ & $(0.998765, 0.0029407)$ & $(0.9987884, 0.0028845)$ \\
		& $\bm P_4$ & $(0.9979064, 0.0050413)$ & $(0.998722, 0.0030827)$ & $(0.9987731, 0.0029601)$ & $(0.9987965, 0.0029039)$ \\ \bottomrule
	\end{tabular}}
\end{table}

\subsection{Test Size}

To analyze the efficiency of the confidence region calculated, we tested its applicability on a set of true random data generated physically not used by the algorithm during its construction. 
We assessed the size of the test by contrasting $100$ new TWNRS for each situation of $D=3,4,5,6$ and of $\alpha=0.01,0.05$.
Table~\ref{tab:result1} and Fig.~\ref{fig:RNG} show the results.

On the one hand, long series ($T=50000$) present a good size for every embedding dimension.
On the other hand, short series ($T=1000$) exhibit only one situation with a noticeable divergence between the expected and the observed size: the test rejects \SI{13}{\percent} of the $100$ series when $D=6$. 
In contrast, we expected \SI{1}{\percent} of rejection.
This might be because, in this case, the condition $D!\ll T$ is not respected.
Notice that the wrongly rejected TWNRS are all close to the point $(1,0)$.

We may then conclude that the test has good empirical size, provided $D!\ll T$, a condition that does not hold for $D=6$ and $T=1000$.

\begin{figure}
    \centering
    \includegraphics[width=\linewidth]{Figures/RNG-1000.pdf}
    \includegraphics[width=\linewidth]{Figures/RNG-50000.pdf}
    \caption{Results of the analysis behavior of true random noises in the regions of confidence built.}
    \label{fig:RNG}
\end{figure}

\subsection{Test Power}

We assessed the power of the test by contrasting time series with different correlation structure (under the $f^{-k}$ model) in the $H \times C$ plane.
Several studies in the literature have used this approach for identifying and characterizing randomness.

Our study's basis is the emblematic time series for each length $T$ and dimension embedding $D$.
Recall that the emblematic time was chosen as the most representative of the data set.
We use these series, transform them into $f^{-k}$ correlated noise, and verify the new point's location in the $H\times C$ plane.

As we can observe in the plane, as the correlation between the observations increases, that is, $k > 0$, the randomness decreases, and the entropy presented decreases, informing the loss of its stochastic characteristic.

Fig.~\ref{fig:CorrNoisea} shows the overall effect of transforming the emblematic time series into $f^{-k}$ correlated noise, with $k=1/2,1,3/2,2,5/2,3$.
At this scale, the emblematic time series $k=0$ and the one with $k=1/2$ appear overlapped.
As the correlation increases with $k$, the randomness decreases, causing a drop in the entropy; the series become progressively more predictable.

Fig.~\ref{fig:CorrNoiseb} is a zoom close to the $(1,0)$ point, along with the confidence regions for the white noise.
We see that $ k = 0 $ and $ k = 0.1 $ are inside the $\SI{95}{\percent}$ confidence region, and $ k = 0.2 $ is inside the $\SI{99}{\percent}$ box.
Notice that the time series with $k=3/10$ is outside the confidence regions and does not pass the randomness test.
The same holds for all $k>3/10$.

\begin{figure*}
	\centering
	\subfigure[Points in the $H\times C$ plane of the emblematic white noise  ($k=0$) and its transformations to become $f^{-k}$ correlated noise with $k=1/2,1,3/2,2,5/2,3$.\label{fig:CorrNoisea}]{
		\includegraphics[width=.48\linewidth]{Figures/Correlation-Analysis-dotted.pdf}}
	\subfigure[Points of the emblematic white noise ($k=0$), and its $f^{-k}$ correlated noise versions, with $k=1/10, 1/5,3/10$ along with the confidence regions for white noise.\label{fig:CorrNoiseb}]{
		\includegraphics[width=.48\linewidth]{Figures/Correlation-Analysis-point.pdf}}
	\caption{Analysis of the test power with correlated $f^{-k}$ noise.}
	\label{fig:correlation}
\end{figure*}

\subsection{Revisiting the White Noise Hypothesis in the Literature}

In this section, we compare the performance of our test with that of 
previous analyses that employ the Complexity-Entropy plane.
To this aim, we produced $100$ sequences of length $T = \num[scientific-notation=true]{5 e4}$ for each generator and computed the $p$-value for each $D = \{3, 4, 5, 6\}$.
Previous results are shown in Table~\ref{Tab:Literature}, and ours are in Table~\ref{Tab:LiteratureComparations}.
We grouped our results in those that rejected (R) the null hypothesis and those that did not reject it (NR).

Comparing Tables~\ref{Tab:Literature} and~\ref{Tab:LiteratureComparations}, we see that our test captures adequately the random dynamics of the sequences produced by most of the analyzed generators.
It is noteworthy that the generator Combo RNG sequences only pass our white noise test for $D = 3$.
In higher embedding dimensions, as we consider longer words, the sequences produced by this generator are not labeled as white noise.

\begin{table}
	\caption{Results of the sequences generated by the main PRNGs in the literature. 
		The sequences have length $T=\num[scientific-notation = true]{5 e4}$.}
	\label{Tab:LiteratureComparations}
	\centering
    \scalebox{0.99}{
	\begin{tabular}{cccc}
		\toprule
		Algorithm & 
		\multicolumn{1}{c}{$D$} & 
		$p$-value &
		HC-PCA\\
		\cmidrule(lr){1-1}
		\cmidrule(lr){2-2}
		\cmidrule(lr){3-3}
		\cmidrule(lr){4-4}
		MOT & 3 & 0.305 & NR\\
		& 4 & 0.572 & NR\\ 
		& 5 & 0.455 & NR\\ 
		& 6 & 0.508 & NR\\ 
		\cmidrule(lr){1-4}
		MWC & 3 & 0.501 & NR\\
		& 4 & 0.477 & NR\\ 
		& 5 & 0.496 & NR\\ 
		& 6 & 0.496 & NR\\ 
		\cmidrule(lr){1-4}
		COM & 3 & 0.123 & NR\\
		& 4 & 0.002 & R\\ 
		& 5 & \num[scientific-notation=true]{1.11 e-16} & R\\ 
		& 6 & \num[scientific-notation=true]{1.11 e-16} & R\\ 
		\cmidrule(lr){1-4}
		LEH & 3 & 0.531 & NR\\
		& 4 & 0.515 & NR\\ 
		\bottomrule
	\end{tabular}
	\begin{tabular}{|cccc}
		\toprule
		Algorithm & 
		\multicolumn{1}{c}{$D$} & 
		$p$-value &
		HC-PCA\\
		\cmidrule(lr){1-1}
		\cmidrule(lr){2-2}
		\cmidrule(lr){3-3}
		\cmidrule(lr){4-4}
		LEH & 5 & 0.495 & NR\\ 
		& 6 & 0.501 & NR\\ 
		\cmidrule(lr){1-4}
		fGn & 3 & 0.521 & NR\\
		& 4 & 0.519 & NR\\ 
		& 5 & 0.498 & NR\\ 
		& 6 & 0.470 & NR\\
		\cmidrule(lr){1-4}
		$f^{-k}$ & 3 & 0.482 & NR\\
		& 4 & 0.520 & NR\\ 
		& 5 & 0.513 & NR\\ 
		& 6 & 0.508 & NR\\
		\cmidrule(lr){1-4}
		LCG & 3 & 0.009 & R\\ 
		& 4 & \num[scientific-notation=true]{1.11 e-16} & R\\ 
		& 5 & \num[scientific-notation=true]{1.11 e-16} & R\\ 
		& 6 & \num[scientific-notation=true]{1.11 e-16} & R\\ 
		\bottomrule
	\end{tabular}}
\end{table}

\section{Conclusions}\label{Sec:Conclusions}

We presented and evaluated the first test for white noise in the  Complexity-Entropy plane.
Our proposal is based on two stages:
(1)~building non-parametric empirical confidence regions in the principal components space and mapping these boxes back to the $H\times C$ plane.
(2)~computing an approximate $p$-value for a given sample by comparing it with points produced by true white noise random sequences (TWNRS).
We obtained the TWNRS with data from physical devices.

Our test has a good size, mostly with long TWNRS.
We also determined the power of our test for the alternative hypothesis of correlated $f^{-k}$ noise and found that it rejects the null hypothesis ($k=0$) for $k>3/10$.

Although our work focuses on the study of short sequences, we were able to capture the random behavior of well-known pseudorandom number generators already analyzed in the literature. 
With this, we verified the adequacy of our technique as it is capable of detecting correlation structures.

\section{Reproducibility and Replicability}\label{Sec:code}

Following the recommendations provided by \cite{ABadgingSystemforReproducibilityandReplicabilityinRemoteSensingResearch}, we make the text, source code, and data used in this study available at the \textit{Confidence-Regions} repository \url{https://github.com/EduardaChagas/ConfidenceRegions}.

\chapter{Conclusions and Future Work}\label{chapter:conclusion}

The main objective of this work was the investigation of problems present in the methodology of symbolization of Bandt-Pompe and its applicability in the characterization of time series and images.
Interested in expanding the range of possible applications, we focus on investigating properties of transition graphs and their possible limitations present in the state of the art.
Another objective was the study of the joint distribution obtained by the descriptors of the Complexity-Entropy plane, as well as possible linear transformations in this space.
Thus, we have advanced in the state of the art by proposing some solutions to deal with scenarios not foreseen in the seminal article by Bandt-Pompe.

The work developed consists of the presentation of two new approaches and their respective applications in data analysis scenarios.
The first problem investigated was the characterization and classification of homogeneous patches of SAR image textures.
Knowing the limitations that exist in sequences that do not follow a unidimensional structure, the first step was the study of linearization methods, as well as the study of the different properties obtained in each class of data analyzed.
Thus, we propose a new approach for classification of remote sensing images, which consists of the following steps: linearization of data using Hilbert curves, generation of the weighted amplitude transition graphs for each sequence, extraction of information theory descriptors, and the classification of these features through the k-nearest neighbors algorithm.
As a result, we verified that we were able to obtain the same evaluation metrics of state-of-the-art handcraft algorithms with the use of only two features, which still provides the user with a new way of the general view of the problem, through the Complexity-Entropy plane.

The most important step of the application above was the proposal for a new generation of the ordinal patterns probability distribution, the weighted amplitude transition graph (WATG).
Considering that the magnitude variation of the targets' backscatter signal is an intrinsic characteristic of the analyzed region, we propose a modification in the ordinal patterns transition graph.
As can be seen in the literature, one of the limitations of traditional methods of ordinal patterns is the absence of amplitude information during the generation of the distribution.
In this way, we propose here the first variation of transition graphs that adds amplitude variation information to its edges.
An advantage of using WATG before classical techniques of adding amplitude information to the symbol histogram is in the higher level of information captured about the underlying dynamics of the system, even when using smaller values of patterns dimension.
Thus, we can obtain a greater degree of data discrimination more quickly and efficiently.

Another focus of the work was the development of an approach to build confidence regions in the Complexity-Entropy plane.
In order to improve the plane characterization capacity, we proposed the HC-PCA, which consists of calculating empirical regions in a new space generated by the linear transformation made by the principal component analysis algorithm.
We found that due to the correlation between the descriptors, proposals such as classical bivariate analysis, regressions and generalized linear models cannot describe the system dynamics well.
To verify the effectiveness of the proposal, we used these regions to capture the randomness of PRNGs in short sequences and thus characterize generators previously analyzed in the literature.
The technique proved to be fast, consistent and robust the addition of correlation structures.

This work presents several possibilities for future research.
For example, the use of WATG can be explored in different application scenarios.
Considering that its main characteristic consists of discriminating sequences with variations in amplitude along with the arrangement of its elements, its applicability is not restricted to remote sensing images.
In the context of SAR images, modifications can be made to increase the generalizability of the technique.
Its possible application in the segmentation of classes of regions is a challenging problem.

On the other hand, under the context of confidence regions, our work opens up a huge range of related research.
The study of regression models on correlated descriptors and the development of specific kernels for the Complexity-Entropy plane are fruitful possibilities for investigation.
We also emphasize the need for efforts to build representative metrics.
With the advancement of deep metric learning techniques~\citep{barros2020new}, we can explore the learning of projections in a linear transformation specific to the plane, which would allow progress to build specific machine learning algorithms for the Complexity-Entropy space.
This means that there is a lot of space to conduct further investigations regarding WATG and HC-PCA.

\ppgccbibliography{bibfile}
% apêndices, se houver
% \begin{appendices}
% \chapter{Um apêndice}
% ...conteúdo do apêndice...
% \end{appendices}
% % anexos, se houver
% \begin{attachments}
% \chapter{Um anexo}
% \end{attachments}
\end{document}